\chapter{Weak Solutions, Part II}

% Source page 0073



In this chapter, we continue to discuss the regularity of weak solutions of elliptic equations of the divergence form. In the first three sections, we discuss the DeGiorgi-Nash-Moser theory for linear elliptic equations of the divergence form. In the first section, we prove the local boundedness of solutions. In the second section, we prove the H\"older continuity. Then in Section 3, we discuss the Harnack inequality. For all results in these three sections, there is no regularity assumption on coefficients.

\section{Local Boundedness}

The main theorem in this section is the following boundedness result.

\begin{theorem}\label{hanlin:weaksol-part2-local-boundedness}\dcref{Theorem 4.1}\group{ch04-weak-solutions-part-ii}\source{han-lin-lecture-notes:page-0073,han-lin-lecture-notes:page-0074,han-lin-lecture-notes:page-0075,han-lin-lecture-notes:page-0076,han-lin-lecture-notes:page-0077,han-lin-lecture-notes:page-0078,han-lin-lecture-notes:page-0079,han-lin-lecture-notes:page-0080,han-lin-lecture-notes:page-0081}
\emph{Suppose $a_{ij} \in L^\infty(B_1)$ and $c \in L^q(B_1)$, for some $q > n/2$, satisfy the following assumptions}

\[
 a_{ij}(x)\xi_i\xi_j \geq \lambda|\xi|^2 \qquad \text{for any } x \in B_1, \xi \in \R^n,
\]

\emph{and}

\[
 |a_{ij}|_{L^\infty} + \|c\|_{L^q} \leq \Lambda,
\]

\emph{for some positive constants $\lambda$ and $\Lambda$. Suppose that $u \in H^1(B_1)$ is a subsolution in the following sense}

\[
(*)\qquad \int_{B_1} a_{ij}D_i u\,D_j\varphi + cu\varphi \leq \int_{B_1} f\varphi
\]
\[
\emph{for any } \varphi \in H^1_0(B_1) \text{ with } \varphi \geq 0 \text{ in } B_1.
\]

\emph{If $f \in L^q(B_1)$, then $u^+ \in L^\infty_{\mathrm{loc}}(B_1)$. Moreover, for any $\theta \in (0,1)$ and any $p > 0$}

\[
\sup_{B_\theta} u^+ \leq C\left\{\frac{1}{(1-\theta)^{n/p}}\|u^+\|_{L^p(B_1)} + \|f\|_{L^q(B_1)}\right\},
\]

\emph{where $C$ is a positive constant depending only on $n,\lambda,\Lambda,p$ and $q$.}
\end{theorem}

In the following, we use two approaches to prove this theorem, one by DeGiorgi and the other by Moser.

\begin{proof}
We first prove for $\theta = 1/2$ and $p = 2$.

\emph{Method 1.} We first present an iteration method due to DeGiorgi.

Consider $v = (u-k)^+$ for $k \geq 0$ and $\zeta \in C^1_0(B_1)$. Set $\varphi = v\zeta^2$ as the test function. Note $v = u-k, Dv = Du$ a.e. in $\{u > k\}$ and $v = 0, Dv = 0$ a.e. in $\{u \leq k\}$. Hence, if we substitute such a $\varphi$ in (*), we integrate in the set $\{u > k\}$.

% Source page 0074

By the Cauchy inequality, we have
\[
\int a_{ij}D_i uD_j\varphi
= \int a_{ij}D_i uD_jv\zeta^2 + 2a_{ij}D_i uD_j\zeta\,v
\]
\[
\ge \lambda \int |Dv|^2\zeta^2 - 2\lambda \int |Dv||D\zeta|v\zeta
\]
\[
\ge \frac{\lambda}{2}\int |Dv|^2\zeta^2 - \frac{2\lambda^2}{\lambda}\int |D\zeta|^2v^2.
\]

Then, we obtain
\[
\int |Dv|^2\zeta^2 \le C\left\{\int v^2|D\zeta|^2 + \int |c|v^2\zeta^2 + k^2\int |c|\zeta^2 + \int |f|v\zeta^2\right\},
\]
and hence
\[
\int |D(v\zeta)|^2 \le C\left\{\int v^2|D\zeta|^2 + \int |c|v^2\zeta^2 + k^2\int |c|\zeta^2 + \int |f|v\zeta^2\right\}.
\]

Recall the Sobolev inequality for $v\zeta\in H_0^1(B_1)$
\[
\left(\int_{B_1}(v\zeta)^{2^*}\right)^{\frac{2}{2^*}} \le c(n)\int_{B_1}|D(v\zeta)|^2,
\]
where $2^* = 2n/(n-2)$ for $n>2$ and $2^*>2$ is arbitrary if $n=2$. The Hölder inequality implies that with $\delta>0$ small and $\zeta\le 1$
\[
\int |f|v\zeta^2 \le \left(\int |f|^q\right)^{\frac{1}{q}}
\left(\int |v\zeta|^{2^*}\right)^{\frac{1}{2^*}}
|\{v\zeta \ne 0\}|^{1-\frac{1}{2^*}-\frac{1}{q}}
\]
\[
\le c(n)\|f\|_{L^q}\left(\int |D(v\zeta)|^2\right)^{\frac{1}{2}}
|\{v\zeta \ne 0\}|^{\frac{1}{2}+\frac{1}{n}-\frac{1}{q}}
\]
\[
\le \delta\int |D(v\zeta)|^2 + c(n,\delta)\|f\|_{L^q}^2|\{v\zeta \ne 0\}|^{1+\frac{2}{n}-\frac{2}{q}}.
\]

Note for $q>n/2$
\[
1+\frac{2}{n}-\frac{2}{q} > 1-\frac{1}{q}.
\]

Therefore, we have
\[
\int |D(v\zeta)|^2 \le C\left\{\int v^2|D\zeta|^2 + \int |c|v^2\zeta^2 + k^2\int |c|\zeta^2 + F^2|\{v\zeta \ne 0\}|^{1-\frac{1}{q}}\right\},
\]
where $F=\|f\|_{L^q(B_1)}$.
We claim
\[
(1)\qquad \int |D(v\zeta)|^2 \le C\left\{\int v^2|D\zeta|^2 + (k^2+F^2)|\{v\zeta \ne 0\}|^{1-\frac{1}{q}}\right\},
\]
if $|\{v\zeta\ne 0\}|$ is small.
% Source page 0075


It is obvious if $c \equiv 0$. In fact, in this special case there is no restriction on the set $\{v\zeta \neq 0\}$. In general, the Hölder inequality implies
\[
\int |c|v^2\zeta^2 \le \left(\int |c|^q\right)^{\frac{1}{q}}
\left(\int (v\zeta)^{2^*}\right)^{\frac{2}{2^*}}
|\{v\zeta \neq 0\}|^{1-\frac{2}{2^*}-\frac{1}{q}}
\]
\[
\le c(n)\int |D(v\zeta)|^2 \left(\int |c|^q\right)^{\frac{1}{q}}
|\{v\zeta \neq 0\}|^{\frac{2}{n}-\frac{1}{q}},
\]
and
\[
\int |c|\zeta^2 \le \left(\int |c|^q\right)^{\frac{1}{q}}|\{v\zeta \neq 0\}|^{1-\frac{1}{q}}.
\]
Therefore, we have
\[
\int |D(v\zeta)|^2 \le C\left\{\int v^2|D\zeta|^2 + \int |D(v\zeta)|^2 \{v\zeta \neq 0\}^{\frac{2}{n}-\frac{1}{q}}
\right.
\]
\[
\left. +\ (k^2 + F^2)|\{v\zeta \neq 0\}|^{1-\frac{1}{q}}\right\}.
\]
This implies (1) if $|\{v\zeta \neq 0\}|$ is small.

To continue, we note by the Sobolev inequality
\[
\int (v\zeta)^2 \le \left(\int (v\zeta)^{2^*}\right)^{\frac{2}{2^*}}|\{v\zeta \neq 0\}|^{1-\frac{2}{2^*}}
\]
\[
\le c(n)\int |D(v\zeta)|^2|\{v\zeta \neq 0\}|^{\frac{2}{n}}.
\]
Therefore, we have
\[
\int (v\zeta)^2 \le C\left\{\int v^2|D\zeta|^2|\{v\zeta \neq 0\}|^{\frac{2}{n}}
 + (k+F)^2|\{v\zeta \neq 0\}|^{1+\frac{2}{n}-\frac{1}{q}}\right\},
\]
if $|\{v\zeta \neq 0\}|$ is small. Hence there exists an $\epsilon > 0$ such that
\[
\int (v\zeta)^2 \le C\left\{\int v^2|D\zeta|^2|\{v\zeta \neq 0\}|^{\epsilon} + (k+F)^2|\{v\zeta \neq 0\}|^{1+\epsilon}\right\},
\]
if $|\{v\zeta \neq 0\}|$ is small. For any fixed $0 < r < R \le 1$, choose $\zeta \in C_0^\infty(B_R)$ such that $\zeta \equiv 1$ in $B_r$, and $0 \le \zeta \le 1$ and $|D\zeta| \le 2(R-r)^{-1}$ in $B_1$. Set
\[
A(k,r) = \{x \in B_r; u \ge k\}.
\]
We conclude that for any $0 < r < R \le 1$ and $k > 0$
\[
\int_{A(k,r)} (u-k)^2 \le C\left\{\frac{1}{(R-r)^2}|A(k,R)|^\epsilon \int_{A(k,R)} (u-k)^2\right.
{}
\left. +\ (k+F)^2|A(k,R)|^{1+\epsilon}\right\},
\tag{2}
\]
if $|A(k,R)|$ is small. Note
\[
|A(k,R)| \le \frac{1}{k}\int_{A(k,R)} u^+ \le \frac{1}{k}\|u^+\|_{L^2}.
\]
Hence (2) holds if $k \ge k_0 = C\|u^+\|_{L^2}$ for some large $C$ depending only on $\lambda$ and $\Lambda$.
% Source page 0076

Next, we claim for $k = C(k_0 + F)$
\[
\int_{A(k,1/2)} (u - k)^2 = 0.
\]
To this end, we take any $h \geq k_0$ and any $0 < r < 1$. It is obvious that
$A(k,r) \supset A(h,r)$. Hence we have
\[
\int_{A(h,r)} (u - h)^2 \leq \int_{A(k,r)} (u - k)^2,
\]
and
\[
|A(h,r)| = |B_r \cap \{u - k > h - k\}| \leq \frac{1}{(h - k)^2}
\int_{A(k,r)} (u - k)^2.
\]
Therefore, by (2) we have for any $h > k \geq k_0$ and $1/2 \leq r < R \leq 1$
\[
\int_{A(h,r)} (u - h)^2
\]
\[
\leq C\left\{ \frac{1}{(R - r)^2} \int_{A(h,R)} (u - h)^2 + (h + F)^2 |A(h,R)| \right\} |A(h,R)|^\varepsilon
\]
\[
\leq C\left\{ \frac{1}{(R - r)^2} + \frac{(h + F)^2}{(h - k)^2} \right\} \frac{1}{(h - k)^{2\varepsilon}}
\left( \int_{A(k,R)} (u - k)^2 \right)^{1+\varepsilon},
\]
or
\[
(3)\qquad \lVert (u - h)^+ \rVert_{L^2(B_r)} \leq C\left\{ \frac{1}{R - r} + \frac{h + F}{h - k} \right\} \frac{1}{(h - k)^\varepsilon}
\lVert (u - k)^+ \rVert_{L^2(B_R)}^{1+\varepsilon}.
\]
Now we carry out the iteration. Set
\[
\varphi(k,r) = \lVert (u - k)^+ \rVert_{L^2(B_r)}.
\]
For $\tau = 1/2$ and some $k > 0$ to be determined, define for $\ell = 0,1,2,\cdots$,
\[
k_\ell = k_0 + k\left(1 - \frac{1}{2^\ell}\right) \qquad (\leq k_0 + k),
\]
\[
r_\ell = \tau + \frac{1}{2^\ell}(1 - \tau).
\]
Obviously, we have
\[
k_\ell - k_{\ell-1} = \frac{k}{2^\ell}, \qquad r_{\ell-1} - r_\ell = \frac{1}{2^\ell}(1 - \tau).
\]
Therefore, we have for $\ell = 0,1,2,\cdots$
\[
\varphi(k_\ell,r_\ell) \leq C\left\{ \frac{2^\ell}{1 - \tau} + \frac{2^\ell(k_0 + F + k)}{k} \right\}
\frac{2^{\ell\varepsilon}}{k^{\varepsilon}} [\varphi(k_{\ell-1},r_{\ell-1})]^{1+\varepsilon}
\]
\[
\leq \frac{C}{1 - \tau} \cdot \frac{k_0 + F + k}{k^{1+\varepsilon}} \cdot 2^{(1+\varepsilon)\ell} \cdot [\varphi(k_{\ell-1},r_{\ell-1})]^{1+\varepsilon}.
\]
Next, we prove inductively for any $\ell = 0,1,\cdots$
\[
(4)\qquad \varphi(k_\ell,r_\ell) \leq \frac{\varphi(k_0,r_0)}{\gamma^\ell}
\qquad \text{for some } \gamma > 1,
\]
if $k$ is sufficiently large. Obviously, it is true for $\ell = 0$. Suppose it is true for $\ell - 1$.
We write
\[
[\varphi(k_{\ell-1},r_{\ell-1})]^{1+\varepsilon} \leq \left\{ \frac{\varphi(k_0,r_0)}{\gamma^{\ell-1}} \right\}^{1+\varepsilon}
= \frac{\varphi(k_0,r_0)^\varepsilon}{\gamma^{\ell\varepsilon-(1+\varepsilon)}} \cdot \frac{\varphi(k_0,r_0)}{\gamma^\ell}.
\]
% Source page 0077
Then, we obtain
\[
\varphi(k\ell,r\ell)\le \frac{C\gamma^{1+\varepsilon}}{1-\tau}\cdot \frac{k_0+F+k}{k^{1+\varepsilon}}\cdot [\varphi(k_0,r_0)]^{\varepsilon}\cdot \frac{2^{\ell(1+\varepsilon)}}{\gamma^{\ell\varepsilon}}\cdot \frac{\varphi(k_0,r_0)}{\gamma^{\ell}}.
\]
Choose $\gamma$ first such that $\gamma^{\varepsilon}=2^{1+\varepsilon}$ and note $\gamma>1$. Next, we need
\[
\frac{C\gamma^{1+\varepsilon}}{1-\tau}\cdot \left(\frac{\varphi(k_0,r_0)}{k}\right)^{\varepsilon}\cdot \frac{k_0+F+k}{k}\le 1.
\]
Therefore, we choose
\[
k=C_{*}\{k_0+F+\varphi(k_0,r_0)\},
\]
for $C_{*}$ large. Let $\ell\to +\infty$ in (4). We conclude
\[
\varphi(k_0+k,\tau)=0.
\]
Hence we have
\[
\sup_{B_{1/2}}u^{+}\le (C_{*}+1)\{k_0+F+\varphi(k_0,r_0)\}.
\]
Recall $k_0=C\|u^{+}\|_{L^{2}(B_1)}$ and $\varphi(k_0,r_0)\le \|u^{+}\|_{L^{2}(B_1)}$. This finishes the proof by the approach due to DeGiorgi.

\emph{Method 2.} We will present an iteration argument due to Moser. First we explain the idea. By choosing a test function appropriately, we estimate $L^{p_1}$ norm of $u$ in a small ball by $L^{p_2}$ norm of $u$ in a large ball for $p_1>p_2$, i.e.,
\[
\|u\|_{L^{p_1}(B_{r_1})}\le C\|u\|_{L^{p_2}(B_{r_2})},
\]
for $p_1>p_2$ and $r_1<r_2$. This is a reversed Hölder inequality. As a sacrifice, $C$ behaves like $\frac{1}{r_2-r_1}$. We then obtain the desired result by an iteration and a careful choice of $\{r_i\}$ and $\{p_i\}$.

For some $k>0$ and $m>0$, set $\bar{u}=u^{+}+k$ and
\[
\bar{u}_m=
\begin{cases}
\bar{u} & \text{if }u<m,\\
k+m & \text{if }u\ge m.
\end{cases}
\]
Then we have $D\bar{u}_m=0$ in $\{u<0\}$ and $\{u>m\}$ and $\bar{u}_m\le \bar{u}$. Consider the test function
\[
\varphi=\eta^{2}(\bar{u}_m^{\beta}\bar{u}-k^{\beta+1})\in H_0^{1}(B_1),
\]
for some $\beta\ge 0$ and some nonnegative function $\eta\in C_0^{1}(B_1)$. A direct calculation yields
\[
D\varphi=\beta\eta^{2}\bar{u}_m^{\beta-1}D\bar{u}_m\bar{u}+\eta^{2}\bar{u}_m^{\beta}D\bar{u}+2\eta D\eta(\bar{u}_m^{\beta}\bar{u}-k^{\beta+1})
\]
\[
=\eta^{2}\bar{u}_m^{\beta}(\beta D\bar{u}_m+D\bar{u})+2\eta D\eta(\bar{u}_m^{\beta}\bar{u}-k^{\beta+1}).
\]
We should emphasize that later on we will begin the iteration with $\beta=0$. Note $\varphi=0$ and $D\varphi=0$ in $\{u\le 0\}$. Hence, if we substitute such a $\varphi$ in (*), we integrate
% Source page 0078

in the set $\{u>0\}$. Note also that $u^+ \le \bar u$ and $\bar u_m^\beta \bar u - k^{\beta+1} \le \bar u_m^\beta \bar u$ for $k > 0$. First,
we have by the Hölder inequality
\[
\int a_{ij}D_i uD_j\varphi
= \int a_{ij}D_i \bar u(BD_j\bar u_m + D_j\bar u)\eta^2\bar u_m^\beta
+ 2\int a_{ij}D_i\bar uD_j\eta(\bar u_m^\beta \bar u - k^{\beta+1})\eta
\]
\[
\ge \lambda\beta\int \eta^2\bar u_m^\beta |D\bar u_m|^2 + \lambda\int \eta^2\bar u_m^\beta |D\bar u|^2 - \Lambda\int |D\bar u||D\eta|\bar u_m^\beta \bar u\eta
\]
\[
\ge \lambda\beta\int \eta^2\bar u_m^\beta |D\bar u_m|^2 + \frac{\lambda}{2}\int \eta^2\bar u_m^\beta |D\bar u|^2 - \frac{2\Lambda^2}{\lambda}\int |D\eta|^2\bar u_m^\beta \bar u^2.
\]
Hence, we obtain by noting $\bar u \ge k$
\[
\beta\int \eta^2\bar u_m^\beta |D\bar u_m|^2 + \int \eta^2\bar u_m^\beta |D\bar u|^2
\]
\[
\le C\left\{\int |D\eta|^2\bar u_m^\beta \bar u^2 + \int \bigl(|c|\eta^2\bar u_m^\beta \bar u^2 + |f|\eta^2\bar u_m^\beta \bar u\bigr)\right\}
\]
\[
\le C\left\{\int |D\eta|^2\bar u_m^\beta \bar u^2 + \int c_0\eta^2\bar u_m^\beta \bar u^2\right\},
\]
where $c_0$ is defined by
\[
c_0 = |c| + \frac{|f|}{k}.
\]

Choose $k = \|f\|_{L^q}$ if $f$ is not identically zero. Otherwise choose an arbitrary $k > 0$
and eventually let $k \to 0+$. By the assumption, we have
\[
\|c_0\|_{L^q} \le \Lambda + 1.
\]

Set $w = \bar u_m^{\beta/2}\bar u$. Note
\[
|Dw|^2 \le (1+\beta)\{\beta\bar u_m^\beta|D\bar u_m|^2 + \bar u_m^\beta|D\bar u|^2\}.
\]

Therefore, we have
\[
\int |Dw|^2\eta^2 \le C\left\{(1+\beta)\int w^2|D\eta|^2 + (1+\beta)\int c_0w^2\eta^2\right\},
\]
or
\[
\int |D(w\eta)|^2 \le C\left\{(1+\beta)\int w^2|D\eta|^2 + (1+\beta)\int c_0w^2\eta^2\right\}.
\]
The Hölder inequality implies
\[
\int c_0w^2\eta^2 \le \left(\int c_0^q\right)^{\frac1q}\left(\int (\eta w)^{\frac{2q}{q-1}}\right)^{1-\frac1q}
\le (\Lambda+1)\left(\int (\eta w)^{\frac{2q}{q-1}}\right)^{1-\frac1q}.
\]
By the interpolation inequality and the Sobolev inequality with
\[
2^* = \frac{2n}{n-2} > \frac{2q}{q-1} > 2 \text{ if } q > n/2,
\]
we have
\[
\|\eta w\|_{L^{\frac{2q}{q-1}}}
\le \epsilon\|w\eta\|_{L^{2^*}} + C(n,q)\epsilon^{-\frac{n}{2q-n}}\|w\|_{L^2}
\]
\[
\le \epsilon\|D(\eta w)\|_{L^2} + C(n,q)\epsilon^{-\frac{n}{2q-n}}\|w\|_{L^2},
\]
% Source page 0079


for any small $\varepsilon>0$. Therefore, we obtain
\[
\int |D(w\eta)|^2 \le C\left\{ (1+\beta)\int w^2|D\eta|^2 + (1+\beta)^{\frac{2q}{2q-n}}\int w^2\eta^2 \right\},
\]
and in particular
\[
\int |D(w\eta)|^2 \le C(1+\beta)^\alpha \int \bigl(|Dw|^2+\eta^2\bigr)w^2,
\]
where $\alpha$ in a positive number depending only on $n$ and $q$. The Sobolev inequality
then implies
\[
\left(\int |\eta w|^{2\chi}\right)^{\frac{1}{\chi}} \le C(1+\beta)^\alpha \int \bigl(|D\eta|^2+\eta^2\bigr)w^2,
\]
where $\chi=\dfrac{n}{n-2}>1$ for $n>2$ and $\chi>2$ for $n=2$. For any $0<r<R\le 1$,
consider an $\eta\in C_0^1(B_R)$ with the property
\[
\eta\equiv 1 \text{ in } B_r \qquad \text{and} \qquad |D\eta|\le \frac{2}{R-r}.
\]
Then, we obtain
\[
\left(\int_{B_r} w^{2\chi}\right)^{\frac{1}{\chi}} \le C\frac{(1+\beta)^\alpha}{(R-r)^2}\int_{B_R} w^2.
\]
Recalling the definition of $w$, we have
\[
\left(\int_{B_r}\bar{u}^{2\chi}\bar{u}_m^{\beta\chi}\right)^{\frac{1}{\chi}}
\le C\frac{(1+\beta)^\alpha}{(R-r)^2}\int_{B_R}\bar{u}^2\bar{u}_m^\beta.
\]
Set $\gamma=\beta+2\ge 2$. Then we obtain
\[
\left(\int_{B_r}\bar{u}_m^{\gamma\chi}\right)^{\frac{1}{\chi}}
\le C\frac{(\gamma-1)^\alpha}{(R-r)^2}\int_{B_R}\bar{u}^\gamma,
\]
provided the integral in the right hand side is bounded. By letting $m\to +\infty$ we
conclude that
\[
\|\bar{u}\|_{L^{\gamma\chi}(B_r)} \le \left(C\frac{(\gamma-1)^\alpha}{(R-r)^2}\right)^{\frac{1}{\gamma}}\|\bar{u}\|_{L^\gamma(B_R)},
\]
provided $\|\bar{u}\|_{L^\gamma(B_R)}<+\infty$, where $C$ is a positive constant depending only on $n,q,\lambda,
\Lambda$, and independent of $\gamma$. The above estimate suggests that we iterate, beginning
with $\gamma=2$, as $2,2\chi,2\chi^2,\cdots$. Now we set for $i=0,1,2,\cdots$,
\[
\gamma_i=2\chi^i \qquad \text{and} \qquad r_i=\frac{1}{2}+\frac{1}{2^{i+1}}.
\]
By $\gamma_i=\chi\gamma_{i-1}$ and $r_{i-1}-r_i=1/2^{i+1}$, we have for $i=1,2,\cdots$,
\[
\|\bar{u}\|_{L^{\gamma_i}(B_{r_i})} \le C(n,q,\lambda,\Lambda)^{\frac{i}{\chi^i}}\|\bar{u}\|_{L^{\gamma_{i-1}}(B_{r_{i-1}})},
\]
provided $\|\bar{u}\|_{L^{\gamma_{i-1}}(B_{r_{i-1}})}<+\infty$. Hence, by an iteration we obtain
\[
\|\bar{u}\|_{L^{\gamma_i}(B_{r_i})} \le C^{\sum \frac{i}{\chi^i}}\|\bar{u}\|_{L^2(B_1)},
\]
and in particular
\[
\left(\int_{B_{1/2}}\bar{u}^{2\chi^i}\right)^{\frac{1}{2\chi^i}}
\le C\left(\int_{B_1}\bar{u}^2\right)^{\frac{1}{2}}.
\]
% Source page 0080




Letting $i \to +\infty$, we get
\[
\sup_{B_{1/2}} \bar u \le C\|\bar u\|_{L^2(B_1)}
\]
or
\[
\sup_{B_{1/2}} u^+ \le C\{\|u^+\|_{L^2(B_1)} + k\}.
\]
Recall the definition of $k$. This finishes the proof for $p = 2$ by the approach due to
Moser.
\end{proof}

\begin{remark}\label{hanlin:weaksol-part2-bounded-test-function}\dcref{Remark 4.2}\group{ch04-weak-solutions-part-ii}\source{han-lin-lecture-notes:page-0080}\uses{hanlin:weaksol-part2-local-boundedness}
If the subsolution $u$ is bounded, we simply take the test function
\[
\varphi = \eta^2(\bar u^{\beta+1} - k^{\beta+1}) \in H^1_0(B_1),
\]
for some $\beta \ge 0$ and some nonnegative function $\eta \in C^1_0(B_1)$.
\end{remark}

\begin{proof}[Proof of Theorem 4.1, continued]
Next, we discuss the general case of Theorem 4.1. This is based on a dilation
argument. Take any $R \le 1$. Define
\[
\tilde u(y) = u(Ry) \qquad \text{for } y \in B_1.
\]
It is easy to see that $\tilde u$ satisfies
\[
\int_{B_1} \tilde a_{ij} D_i \tilde u D_j \varphi + \tilde c\,\tilde u\varphi \le \int_{B_1} \tilde f\,\varphi
\]
for any $\varphi \in H^1_0(B_1)$ and $\varphi \ge 0$ in $B_1$,
where
\[
\tilde a(y) = a(Ry), \qquad \tilde c(y) = R^2 c(Ry), \qquad \tilde f(y) = R^2 f(Ry),
\]
for any $y \in B_1$. A direct calculation shows
\[
\|\tilde a_{ij}\|_{L^\infty(B_1)} + \|\tilde c\|_{L^q(B_1)} = \|a_{ij}\|_{L^\infty(B_R)} + R^{2-\frac{n}{q}}\|c\|_{L^q(B_R)} \le \Lambda.
\]

We may apply what we just proved to $\tilde u$ in $B_1$ and rewrite the result in terms of $u$.
Hence, we obtain for $p \ge 2$
\[
\sup_{B_{R/2}} u^+ \le C\left\{\frac{1}{R^{n/p}}\|u^+\|_{L^p(B_R)} + R^{2-\frac{n}{q}}\|f\|_{L^q(B_R)}\right\},
\]
where $C$ is a positive constant depending only on $n$, $\lambda$, $\Lambda$, $p$ and $q$.
The estimate in $B_{\theta R}$ can be obtained by applying the above result to $B_{(1-\theta)R}(y)$
for any $y \in B_{\theta R}$. By taking $R = 1$, we have Theorem 4.1 for any $\theta \in (0,1)$ and
$p \ge 2$.

Now we prove the statement for $p \in (0,2)$. We showed that there holds for any
$\theta \in (0,1)$ and $0 < R \le 1$
\[
\|u^+\|_{L^\infty(B_{\theta R})} \le C\left\{\frac{1}{[(1-\theta)R]^{n/2}}\|u^+\|_{L^2(B_R)} + R^{2-\frac{n}{q}}\|f\|_{L^q(B_R)}\right\}
\]
\[
\le C\left\{\frac{1}{[(1-\theta)R]^{n/2}}\|u^+\|_{L^2(B_R)} + \|f\|_{L^q(B_1)}\right\}.
\]
For $p \in (0,2)$, we have
\[
\int_{B_R} (u^+)^2 \le \|u^+\|_{L^\infty(B_R)}^{2-p}\int_{B_R}(u^+)^p,
\]

% Source page 0081



and hence by the Hölder inequality
\[
\|u^+\|_{L^\infty(B_{\theta R})}
\le C\left\{\frac{1}{[(1-\theta)R]^{\frac{n}{2}}}\|u^+\|_{L^\infty(B_R)}^{1-\frac{2}{p}}
\left(\int_{B_R} (u^+)^p \ud x\right)^{\frac{1}{2}} + \|f\|_{L^q(B_R)}\right\}
\]
\[
\le \frac{1}{2}\|u^+\|_{L^\infty(B_R)} + C\left\{\frac{1}{[(1-\theta)R]^{\frac{n}{p}}}\left(\int_{B_R} (u^+)^p\right)^{\frac{1}{p}} + \|f\|_{L^q(B_R)}\right\}.
\]

Set $f(t)=\|u^+\|_{L^\infty(B_t)}$ for $t\in(0,1]$. Then, we obtain for any $0<r<R\le 1$
\[
f(r)\le \frac{1}{2}f(R)+\frac{C}{(R-r)^{\frac{n}{p}}}\|u^+\|_{L^p(B_1)}+C\|f\|_{L^q(B_1)}.
\]

We apply the following lemma to get for any $0<r<R<1$
\[
f(r)\le \frac{C}{(R-r)^{\frac{n}{p}}}\|u^+\|_{L^p(B_1)}+C\|f\|_{L^q(B_1)}.
\]

By letting $R\to 1$, we get for any $\theta<1$
\[
\|u^+\|_{L^\infty(B_\theta)}\le \frac{C}{(1-\theta)^{\frac{n}{p}}}\|u^+\|_{L^p(B_1)}+C\|f\|_{L^q(B_1)}.
\]

This finishes the proof.
\end{proof}

We need the following simple lemma.

\begin{lemma}\label{hanlin:weaksol-part2-iteration-lemma}\dcref{Lemma 4.3}\group{ch04-weak-solutions-part-ii}\source{han-lin-lecture-notes:page-0081}
\emph{Let $f(t)\ge 0$ be bounded in $[\tau_0,\tau_1]$ with $\tau_0\ge 0$. Suppose for any $\tau_0\le t<s\le \tau_1$}
\[
f(t)\le \theta f(s)+\frac{A}{(s-t)^\alpha}+B,
\]
\emph{for some $\theta\in[0,1)$. Then there holds for any $\tau_0\le t<s\le \tau_1$}
\[
f(t)\le c\left\{\frac{A}{(s-t)^\alpha}+B\right\},
\]
\emph{where $c$ is a positive constant depending only on $\alpha$ and $\theta$.}
\end{lemma}

\begin{proof}
Fix $\tau_0\le t<s\le \tau_1$. For some $0<\tau<1$, we consider the sequence $\{t_i\}$ defined by
\[
t_0=t \text{ and } t_{i+1}=t_i+(1-\tau)\tau^i(s-t).
\]
Note $t_\infty=s$. By an iteration, we get
\[
f(t)=f(t_0)\le \theta^k f(t_k)+\left[\frac{A}{(1-\tau)^\alpha}(s-t)^{-\alpha}+B\right]\sum_{i=0}^{k-1}\theta^i\tau^{-i\alpha}.
\]
Choose $\tau<1$ such that $\theta\tau^{-\alpha}<1$, i.e., $\theta<\tau^\alpha<1$. As $k\to+\infty$, we have
\[
f(t)\le c(\alpha,\theta)\left\{\frac{A}{(1-\tau)^\alpha}(s-t)^{-\alpha}+B\right\}.
\]

This finishes the proof.
\end{proof}

In the rest of this section, we use the Moser's iteration to prove a high integrability result, which is closely related to Theorem 4.1. For the next result, we require $n\ge 3$.

% Source page 0082

\begin{theorem}\label{hanlin:weaksol-part2-high-integrability}\dcref{Theorem 4.4}\group{ch04-weak-solutions-part-ii}\source{han-lin-lecture-notes:page-0082,han-lin-lecture-notes:page-0083}\uses{hanlin:weaksol-part2-local-boundedness}
Suppose $a_{ij} \in L^\infty(B_1)$ and $c \in L^{n/2}(B_1)$ satisfy the following assumption
\[
\lambda |\xi|^2 \le a_{ij}(x)\xi_i\xi_j \le \Lambda |\xi|^2
\qquad \text{for any $x \in B_1, \xi \in \mathbb{R}^n$,}
\]
for some positive constants $\lambda$ and $\Lambda$. Suppose that $u \in H^1(B_1)$ is a subsolution in the following sense
\[
\int_{B_1} a_{ij}D_i u D_j \varphi + cu\varphi \le \int_{B_1} f\varphi
\]
for any $\varphi \in H_0^1(B_1)$ and $\varphi \ge 0$ in $B_1$.

If $f \in L^q(B_1)$ for some $q \in \left[\frac{2n}{n+2}, \frac{n}{2}\right)$, then $u^+ \in L^{q^*}_{\mathrm{loc}}(B_1)$ for
\[
\frac{1}{q^*} = \frac{1}{q} - \frac{2}{n}.
\]
Moreover,
\[
\|u^+\|_{L^{q^*}(B_1)} \le C\bigl\{\|u^+\|_{L^2(B_1)} + \|f\|_{L^q(B_1)}\bigr\},
\]
where $C$ is a positive constant depending only on $n, \lambda, \Lambda, q$ and $\varepsilon(K)$, with
\[
\varepsilon(K) = \left(\int_{\{|c|>K\}} |c|^{n/2}\right)^{2/n}.
\]
\end{theorem}

\begin{proof}
For $m > 0$, set $\bar{u} = u^+$ and
\[
\bar{u}_m =
\begin{cases}
\bar{u} & \text{if } u < m,\\
m & \text{if } u \ge m.
\end{cases}
\]
Then consider a test function
\[
\varphi = \eta^2 \bar{u}^{\beta}\bar{u}_m \in H_0^1(B_1),
\]
for some $\beta \ge 0$ and some nonnegative function $\eta \in C_0^1(B_1)$. By similar calculations
as in the proof of Theorem 4.1, we obtain
\[
\left(\int \eta^{2\chi} \bar{u}_m^{\beta\chi} u^{2\chi}\right)^{1/\chi}
\le C(1+\beta)\left\{\int |D\eta|^2 \bar{u}_m^\beta \bar{u}^2 + \int |c|\eta^2 \bar{u}_m^\beta \bar{u}^2 + \int |f|\eta^2 \bar{u}_m^\beta \bar{u}\right\},
\]
where $\chi = \frac{n}{n-2} > 1$. The Hölder inequality implies for any $K > 0$
\[
\int |c|\eta^2 \bar{u}_m^\beta \bar{u}^2 \le K \int_{\{|c|\le K\}} \eta^2 \bar{u}_m^\beta \bar{u}^2 + \int_{\{|c|>K\}} |c|\eta^2 \bar{u}_m^\beta \bar{u}^2
\]
\[
\le K \int \eta^2 \bar{u}_m^\beta \bar{u}^2 + \left(\int_{\{|c|>K\}} |c|^{n/2}\right)^{2/n}\left(\int (\eta^2 \bar{u}_m^\beta \bar{u}^2)^{\frac{n}{n-2}}\right)^{\frac{n-2}{n}}
\]
\[
\le K \int \eta^2 \bar{u}_m^\beta \bar{u}^2 + \varepsilon(K)\left(\int \eta^{2\chi}\bar{u}_m^{\beta\chi}u^{2\chi}\right)^{1/\chi}.
\]
Note $\varepsilon(K) \to 0$ as $K \to +\infty$, since $c \in L^{n/2}(B_1)$. Hence, for bounded $\beta$ we obtain
by choosing large $K = K(\beta)$
\[
\left(\int \eta^{2\chi}\bar{u}_m^{\beta\chi}u^{2\chi}\right)^{1/\chi}
\le C(1+\beta)\left\{\int \bigl(|D\eta|^2 + \eta^2\bigr)\bar{u}_m^\beta \bar{u}^2 + \int |f|\eta^2 \bar{u}_m^\beta \bar{u}\right\}.
\]
Note
\[
\bar u_m^\beta \bar u
\le \bar u_m^{\beta-\frac{\beta}{\beta+2}}\bar u^{1+\frac{\beta}{\beta+2}}
= (\bar u_m^\beta \bar u^2)^{\frac{\beta+1}{\beta+2}}.
\]
% Source page 0083




Therefore, by the H\"older inequality again we have for $\eta \le 1$
\[
\int |f|\eta^2 \bar u_m^\beta \bar u
\le \left(\int |f|^q\right)^{\frac{1}{q}}
\left(\int (\eta^2 \bar u_m^\beta \bar u^2)^\chi\right)^{\frac{\beta+1}{(\beta+2)\chi}}
\|\eta\|_{L^\infty}^{1-\frac{1}{q}-\frac{\beta+1}{(\beta+2)\chi}}
\]
\[
\le \epsilon \left(\int \eta^{2\chi}\bar u_m^\beta \bar u^{2\chi}\right)^{\frac{1}{\chi}}
C(\epsilon,\beta)\left(\int |f|^q\right)^{\frac{\beta+2}{q}},
\]
provided
\[
1-\frac{1}{q}-\frac{\beta+1}{(\beta+2)\chi}\ge 0,
\]
which is equivalent to
\[
\beta+2\le \frac{q(n-2)}{n-2q}.
\]
Hence, $\beta$ is required to be bounded, depending only on $n$ and $q$. Then we obtain
\[
\left(\int \eta^{2\chi}\bar u_m^\beta \bar u^{2\chi}\right)^{\frac{1}{\chi}}
\le C\left\{\int \left(|D\eta|^2+\eta^2\right)\bar u_m^\beta \bar u^2+\|f\|_{L^q}^{\beta+2}\right\}.
\]
By setting $\gamma=\beta+2$, we have by the definition of $q^*$
\[
(1)\qquad 2\le \gamma \le \frac{q(n-2)}{n-2q}=\frac{q^*}{\chi}.
\]
We conclude, as before, for any such $\gamma$ in (1) and any $0<r<R\le 1$
\[
(2)\qquad \|\bar u\|_{L^{\chi\gamma}(B_r)}\le C\left\{\frac{1}{(R-r)^{\frac{2}{\gamma}}}\|\bar u\|_{L^\gamma(B_R)}+\|f\|_{L^q(B_1)}\right\},
\]
provided $\|\bar u\|_{L^\gamma(B_R)}<+\infty$. Again this suggests an iteration $2,2\chi,2\chi^2,\cdots$.

For a given $q\in\left[\frac{2n}{n+2},\frac{n}{2}\right)$, there exists a positive integer $k$ such that
\[
2\chi^{k-1}\le \frac{q(n-2)}{n-2q}<2\chi^k.
\]
For such a $k$, we get by finitely many iterations of (2)
\[
\|\bar u\|_{L^{2\chi^k}(B_{3/4})}\le C\left\{\|\bar u\|_{L^2(B_1)}+\|f\|_{L^q(B_1)}\right\},
\]
and in particular
\[
\|\bar u\|_{L^{q^*}_\chi(B_{3/4})}\le C\left\{\|\bar u\|_{L^2(B_1)}+\|f\|_{L^q(B_1)}\right\}.
\]
With $\gamma=\frac{q^*}{\chi}$ in (2), we obtain
\[
\|\bar u\|_{L^{q^*}(B_{1/2})}\le C\left\{\|\bar u\|_{L^{q^*}_\chi(B_{3/4})}+\|f\|_{L^q(B_1)}\right\}.
\]
This finishes the proof.
\end{proof}

% Source page 0084


\section{The H\"older Continuity}

We first discuss homogeneous equations with no lower order terms. Consider
\[
Lu \equiv -D_i(a_{ij}(x)D_j u) \qquad \text{in } B_1 \subset \mathbb{R}^n,
\]
where $a_{ij} \in L^\infty(B_1)$ satisfies
\[
\lambda |\xi|^2 \le a_{ij}(x)\xi_i\xi_j \le \Lambda |\xi|^2 \qquad \text{for any } x \in B_1 \text{ and } \xi \in \mathbb{R}^n,
\]
for some positive constants $\lambda$ and $\Lambda$.

\begin{definition}\label{hanlin:weaksol-part2-homogeneous-subsolution}\dcref{Definition 4.5}\group{ch04-weak-solutions-part-ii}\source{han-lin-lecture-notes:page-0084}
The function $u \in H^1_{\mathrm{loc}}(B_1)$ is called a subsolution (supersolution) of the equation $Lu = 0$, if
\[
\int_{B_1} a_{ij} D_i u D_j \varphi \le 0 \ (\ge 0),
\]
for any $\varphi \in H^1_0(B_1)$ and $\varphi \ge 0$.
\end{definition}

\begin{lemma}\label{hanlin:weaksol-part2-convex-composition}\dcref{Lemma 4.6}\group{ch04-weak-solutions-part-ii}\source{han-lin-lecture-notes:page-0084}\uses{hanlin:weaksol-part2-homogeneous-subsolution}
Let $\Phi \in C^1_{\mathrm{loc}}(\mathbb{R})$ be convex.
\begin{enumerate}
\item[(i)] If $u$ is a subsolution and $\Phi' \ge 0$, then $v=\Phi(u)$ is a subsolution provided $v \in H^1_{\mathrm{loc}}(B_1)$;
\item[(ii)] If u is a supersolution and $\Phi' \le 0$, then $v=\Phi(u)$ is a subsolution provided $v \in H^1_{\mathrm{loc}}(B_1)$.
\end{enumerate}
\end{lemma}

\begin{remark}\label{hanlin:weaksol-part2-positive-part-subsolution}\dcref{Remark 4.7}\group{ch04-weak-solutions-part-ii}\source{han-lin-lecture-notes:page-0084}\uses{hanlin:weaksol-part2-convex-composition}
If $u$ is a subsolution, then $(u-k)^+$ is also a subsolution, where
\[
(u-k)^+ = \max\{0,u-k\}.
\]
In this case, $\Phi(s) = (s-k)^+$.
\end{remark}

\begin{proof}[Proof of Lemma 4.6]
We only prove (i). Assume first $\Phi \in C^2_{\mathrm{loc}}(\mathbb{R})$. Then we have $\Phi'(s)\ge 0$ and $\Phi''(s)\ge 0$. Consider a $\varphi \in C_0^1(B_1)$ with $\varphi \ge 0$. A direct calculation yields
\[
\int_{B_1} a_{ij} D_i v D_j \varphi
= \int_{B_1} a_{ij}\Phi'(u)D_i u D_j \varphi
\]
\[
= \int_{B_1} a_{ij} D_i u D_j(\Phi'(u)\varphi) - \int_{B_1} (a_{ij}D_i u D_j u)\varphi \Phi''(u) \le 0,
\]
where $\Phi'(u)\varphi \in H_0^1(B_1)$ is nonnegative. In general, set $\Phi_\varepsilon(s)=\rho_\varepsilon * \Phi(s)$ with $\rho_\varepsilon$ as the standard mollifier. Then $\Phi_\varepsilon'(s)=\rho_\varepsilon * \Phi'(s)\ge 0$ and $\Phi_\varepsilon''(s)\ge 0$. Hence $\Phi_\varepsilon(u)$ is a subsolution by what we just proved. Note $\Phi_\varepsilon'(s)\to \Phi'(s)$ a.e. as $\varepsilon \to 0^+$. The Lebesgue dominant convergence theorem implies the desired result.
\end{proof}

We also need the following Poincar\'e-Sobolev inequality.

\begin{lemma}\label{hanlin:weaksol-part2-poincare-sobolev-zero-set}\dcref{Lemma 4.8}\group{ch04-weak-solutions-part-ii}\source{han-lin-lecture-notes:page-0084,han-lin-lecture-notes:page-0085}
Let $\epsilon$ be a positive constant. For any $u \in H^1(B_1)$, if
\[
|\{x \in B_1; u=0\}| \ge \epsilon |B_1|,
\]
then
\[
\int_{B_1} u^2 \le C \int_{B_1} |Du|^2,
\]
where $C$ is a positive constant depending only on $\epsilon$ and $n$.
\end{lemma}

\begin{proof}
We prove by contradiction. If not, there exists a sequence $\{u_m\} \subset H^1(B_1)$ such that
\[
|\{x \in B_1; u_m=0\}| \ge \epsilon |B_1|,
\]
% Source page 0085

and
\[
\int_{B_1} u_m^2 = 1, \qquad \int_{B_1} |D u_m|^2 \to 0 \text{ as } m \to \infty.
\]
We may assume $u_m \to u_0$ strongly in $L^2(B_1)$ and weakly in $H^1(B_1)$, with $u_0 \in H^1(B_1)$. Clearly $u_0$ is a nonzero constant. Then we have
\[
0 = \lim_{m\to\infty} \int_{B_1} |u_m-u_0|^2 \ge \lim_{m\to\infty} \int_{\{u_m=0\}} |u_m-u_0|^2
\]
\[
\ge |u_0|^2 \inf_m |\{u_m=0\}| > 0.
\]
This is a contradiction.
\end{proof}

Now, we begin to discuss the Hölder continuity. We first prove the following result, which is often referred to as the density theorem.

\begin{theorem}\label{hanlin:weaksol-part2-density-theorem}\dcref{Theorem 4.9}\group{ch04-weak-solutions-part-ii}\source{han-lin-lecture-notes:page-0085,han-lin-lecture-notes:page-0086}\uses{hanlin:weaksol-part2-local-boundedness,hanlin:weaksol-part2-convex-composition,hanlin:weaksol-part2-poincare-sobolev-zero-set}
\emph{Suppose $\epsilon \in (0,1)$ is a constant and $u$ is a positive supersolution in $B_2$ with}
\[
|\{x \in B_1; u \ge 1\}| \ge \epsilon |B_1|.
\]
\emph{Then}
\[
\inf_{B_{1/2}} u \ge C,
\]
\emph{where $C$ is a positive constant depending only on $\epsilon$, $n$ and $\Lambda/\lambda$.}
\end{theorem}

\begin{proof}
We may assume $u \ge \delta > 0$ and then let $\delta \to 0^+$ at the end.
By Lemma 4.6, $v = (\log u)^-$ is a subsolution, bounded by $\log \delta^{-1}$. Theorem 4.1 yields
\[
\sup_{B_{1/2}} v \le C \left(\int_{B_1} |v|^2\right)^{\frac12}.
\]
Note $|\{x \in B_1; v = 0\}| = |\{x \in B_1; u \ge 1\}| \ge \epsilon |B_1|$. Lemma 4.8 implies
\begin{equation}
\sup_{B_{1/2}} v \le C \left(\int_{B_1} |Dv|^2\right)^{\frac12}.
\tag{1}
\end{equation}
We will prove that the right-hand side is bounded. To this end, consider a test function $\varphi = \frac{\zeta^2}{u}$ for $\zeta \in C_0^1(B_2)$. Then we obtain
\[
0 \le \int a_{ij} D_i u D_j\!\left(\frac{\zeta^2}{u}\right)
= -\int \zeta^2 \frac{a_{ij}D_i u D_j u}{u^2} + 2\int \frac{\zeta a_{ij}D_i u D_j \zeta}{u},
\]
and hence
\[
\int \zeta^2 |D \log u|^2 \le C \int |D\zeta|^2.
\]
So for fixed $\zeta \in C_0^1(B_2)$ with $\zeta \equiv 1$ in $B_1$, we have
\[
\int_{B_1} |D \log u|^2 \le C.
\]
Combining with (1), we obtain
\[
\sup_{B_{1/2}} v = \sup_{B_{1/2}} (\log u)^- \le C,
\]
% Source page 0086



or

\[
\inf_{B_{1/2}} u \ge e^{-C} > 0.
\]

This finishes the proof.
\end{proof}

The next result controls the oscillation of solutions.

\begin{theorem}\label{hanlin:weaksol-part2-oscillation-decay}\dcref{Theorem 4.10}\group{ch04-weak-solutions-part-ii}\source{han-lin-lecture-notes:page-0086}\uses{hanlin:weaksol-part2-density-theorem}
\emph{Suppose $u$ is a bounded solution of $Lu = 0$ in $B_2$. Then}
\[
\operatorname{osc}_{B_{1/2}} u \le \gamma \operatorname{osc}_{B_1} u,
\]
\emph{where $\gamma \in (0,1)$ is a constant depending only on $n$ and $\Lambda/\lambda$.}
\end{theorem}

\begin{proof}
The local boundedness is proved in the previous section. Set
\[
\alpha_1 = \sup_{B_1} u \qquad \text{and} \qquad \beta_1 = \inf_{B_1} u.
\]
Consider
\[
\frac{u - \beta_1}{\alpha_1 - \beta_1}
\qquad \text{and} \qquad
\frac{\alpha_1 - u}{\alpha_1 - \beta_1}.
\]
Obviously, they are solutions of $Lv = 0$. Note the following equivalence
\[
u \ge \frac{1}{2}(\alpha_1 + \beta_1) \Longleftrightarrow \frac{u - \beta_1}{\alpha_1 - \beta_1} \ge \frac{1}{2},
\]
\[
u \le \frac{1}{2}(\alpha_1 + \beta_1) \Longleftrightarrow \frac{\alpha_1 - u}{\alpha_1 - \beta_1} \ge \frac{1}{2}.
\]

\emph{Case 1.} Suppose
\[
\bigl|\{x \in B_1; \frac{2(u - \beta_1)}{\alpha_1 - \beta_1} \ge 1\}\bigr| \ge \frac{1}{2}|B_1|.
\]
Apply Theorem 4.9 to $\dfrac{u - \beta_1}{\alpha_1 - \beta_1} \ge 0$ in $B_1$. We have for some $C > 1$
\[
\inf_{B_{1/2}} \frac{u - \beta_1}{\alpha_1 - \beta_1} \ge \frac{1}{C},
\]
and hence
\[
\inf_{B_{1/2}} u \ge \beta_1 + \frac{1}{C}(\alpha_1 - \beta_1).
\]

\emph{Case 2.} Suppose
\[
\bigl|\{x \in B_1; \frac{2(\alpha_1 - u)}{\alpha_1 - \beta_1} \ge 1\}\bigr| \ge \frac{1}{2}|B_1|.
\]
Similarly, we obtain
\[
\sup_{B_{1/2}} u \le \alpha_1 - \frac{1}{C}(\alpha_1 - \beta_1).
\]

Now, we set
\[
\alpha_2 = \sup_{B_{1/2}} u \qquad \text{and} \qquad \beta_2 = \inf_{B_{1/2}} u.
\]
Note that $\beta_2 \ge \beta_1$ and $\alpha_2 \le \alpha_1$. In both cases, we have
\[
\alpha_2 - \beta_2 \le \left(1 - \frac{1}{C}\right)(\alpha_1 - \beta_1).
\]

This finishes the proof.
\end{proof}

% Source page 0087

The DeGiorgi theorem on Hölder continuity is an easy consequence of above
results.

\begin{theorem}\label{hanlin:weaksol-part2-degiorgi-holder}\dcref{Theorem 4.11}\group{ch04-weak-solutions-part-ii}\source{han-lin-lecture-notes:page-0087}\uses{hanlin:weaksol-part2-local-boundedness,hanlin:weaksol-part2-oscillation-decay}
\emph{Suppose $u$ is an $H^1(B_1)$ solution of $Lu = 0$ in $B_1$. Then
$u \in C^\alpha(B_1)$ for some $\alpha \in (0,1)$ depending only on $n$ and
$\Lambda/\lambda$. Moreover,}
\[
\sup_{B_{1/2}} |u(x)| + \sup_{x,y \in B_{1/2}} \frac{|u(x) - u(y)|}{|x-y|^\alpha}
\le C\|u\|_{L^2(B_1)},
\]
\emph{where $C$ is a positive constant depending only on $n$ and $\Lambda/\lambda$.}
\end{theorem}

In the rest of this section, we discuss the Hölder continuity of solutions of
general linear equations. We need the following lemma.

\begin{lemma}\label{hanlin:weaksol-part2-energy-decay}\dcref{Lemma 4.12}\group{ch04-weak-solutions-part-ii}\source{han-lin-lecture-notes:page-0087,han-lin-lecture-notes:page-0088}\uses{hanlin:weaksol-part2-degiorgi-holder}
\emph{Suppose that $a_{ij} \in L^\infty(B_r)$ satisfies}
\[
\lambda |\xi|^2 \le a_{ij}(x)\xi_i \xi_j \le \Lambda |\xi|^2 \qquad \text{for any $x \in B_r$, $\xi \in \mathbb{R}^n$},
\]
\emph{for some $0 < \lambda \le \Lambda < +\infty$. Suppose $u \in H^1(B_r)$ satisfies}
\[
\int_{B_r} a_{ij} D_i u D_j \varphi = 0 \qquad \text{for any $\varphi \in H_0^1(B_r)$}.
\]
\emph{Then there exists an $\alpha \in (0,1)$ such that for any $0 < \rho < r$}
\[
\int_{B_\rho} |Du|^2 \le C \left(\frac{\rho}{r}\right)^{n-2+2\alpha} \int_{B_r} |Du|^2,
\]
\emph{where $C$ and $\alpha$ depend only on $n$ and $\Lambda/\lambda$.}
\end{lemma}

\begin{proof}
By a dilation, we consider $r = 1$. We restrict our consideration to the
range $\rho \in (0, \frac14]$, since it is trivial for $\rho \in (\frac14,1]$. We may further assume
$\int_{B_1} u = 0$, since the function $u-|B_1|^{-1}\int_{B_1}u$ solves the same equation. The Poincaré inequality
yields
\[
\int_{B_1} u^2 \le c(n)\int_{B_1} |Du|^2.
\]

Hence Theorem 4.11 implies for $|x| \le 1/2$
\[
|u(x) - u(0)|^2 \le C|x|^{2\alpha}\int_{B_1} |Du|^2,
\]
where $\alpha \in (0,1)$ is as determined in Theorem 4.11. For any $0 < \rho \le 1/4$, take a
cut-off function $\zeta \in C_0^\infty(B_{2\rho})$ with $\zeta \equiv 1$ in $B_\rho$, and $0 \le \zeta \le 1$ and $|D\zeta| \le 2/\rho$.
Then set $\varphi = \zeta^2(u - u(0))$. Now the equation yields
\[
0 = \int_{B_1} a_{ij} D_i u \bigl(\zeta^2 D_j u + 2\zeta D_j\zeta (u-u(0))\bigr)
\]
\[
\ge \frac{\lambda}{2}\int_{B_{2\rho}} \zeta^2 |Du|^2 - C \sup_{B_{2\rho}} |u-u(0)|^2 \int_{B_{2\rho}} |D\zeta|^2.
\]

Therefore, we have
\[
\int_{B_\rho} |Du|^2 \le C\rho^{n-2} \sup_{B_{2\rho}} |u-u(0)|^2.
\]

The conclusion follows easily.
\end{proof}
% Source page 0088



Now, we prove the following result in the same way we proved Theorem 3.8 in
Chapter 3, with Lemma 3.9 in Chapter 3 replaced by Lemma 4.12.

\begin{theorem}\label{hanlin:weaksol-part2-general-linear-holder}\dcref{Theorem 4.13}\group{ch04-weak-solutions-part-ii}\source{han-lin-lecture-notes:page-0088}\uses{hanlin:weaksol-part1-holder-sol,hanlin:weaksol-part2-energy-decay}
\emph{Assume $a_{ij} \in L^\infty(B_1)$ and $c \in L^n(B_1)$ satisfy}
\[
\lambda |\xi|^2 \le a_{ij}(x)\xi_i\xi_j \le \Lambda |\xi|^2 \qquad \text{for any $x \in B_1$, $\xi \in \mathbb{R}^n$,}
\]
\emph{for some $0 < \lambda \le \Lambda < +\infty$. Suppose that $u \in H^1(B_1)$ satisfies}
\[
\int_{B_1} a_{ij}D_j uD_i\varphi + cu\varphi = \int_{B_1} f\varphi \qquad \text{for any $\varphi \in H^1_0(B_1)$.}
\]
\emph{If $f \in L^q(B_1)$ for some $q > n/2$, then $u \in C^\alpha(B_1)$ for some $\alpha \in (0,1)$, depending
only on $n$, $q$, $\lambda$ and $\Lambda$. Moreover, there exists an $R_0 \in (0,1)$ such that for any
$x \in B_{1/2}$ and $r \le R_0$}
\[
\int_{B_r(x)} |Du|^2 \le Cr^{n-2+2\alpha}\left\{\|f\|_{L^q(B_1)}^2 + \|u\|_{H^1(B_1)}^2\right\},
\]
\emph{where $R_0$ and $C$ are constants depending only on $n$, $q$, $\lambda$, $\Lambda$ and $\|c\|_{L^n}$.}
\end{theorem}

\bigskip

\section{Moser's Harnack Inequality}

In this section, we only discuss equations without lower order terms. Suppose
$\Omega$ is a domain in $\mathbb{R}^n$. We always assume that $a_{ij} \in L^\infty(\Omega)$ satisfies
\[
\lambda |\xi|^2 \le a_{ij}(x)\xi_i\xi_j \le \Lambda |\xi|^2 \qquad \text{for any $x \in \Omega$ and $\xi \in \mathbb{R}^n$,}
\]
for some positive constants $\lambda$ and $\Lambda$.

We first state a result of local boundedness, which is simply a dilated version
of Theorem 4.1.

\begin{theorem}\label{hanlin:weaksol-part2-scaled-local-boundedness}\dcref{Theorem 4.14}\group{ch04-weak-solutions-part-ii}\source{han-lin-lecture-notes:page-0088}\uses{hanlin:weaksol-part2-local-boundedness}
\emph{Let $u \in H^1(\Omega)$ be a nonnegative subsolution in $\Omega$ in the following sense}
\[
\int_\Omega a_{ij}D_i uD_j\varphi \le \int_\Omega f\varphi \qquad \text{for any $\varphi \in H^1_0(\Omega)$ with $\varphi \ge 0$ in $\Omega$.}
\]
\emph{Suppose $f \in L^q(\Omega)$ for some $q > n/2$. Then for any $B_R \subset \Omega$, any $0 < r < R$ and any $p > 0$}
\[
\sup_{B_r} u \le C\left(\frac{1}{(R-r)^{n/p}}\|u^+\|_{L^p(B_R)} + R^{2-\frac{n}{q}}\|f\|_{L^q(B_R)}\right),
\]
\emph{where $C$ is a positive constant depending only on $n$, $\lambda$, $\Lambda$, $p$ and $q$.}
\end{theorem}

The next result is referred to as the weak Harnack inequality.

\begin{theorem}\label{hanlin:weaksol-part2-weak-harnack}\dcref{Theorem 4.15}\group{ch04-weak-solutions-part-ii}\source{han-lin-lecture-notes:page-0088,han-lin-lecture-notes:page-0089,han-lin-lecture-notes:page-0090,han-lin-lecture-notes:page-0091,han-lin-lecture-notes:page-0092,han-lin-lecture-notes:page-0093,han-lin-lecture-notes:page-0094}\uses{hanlin:weaksol-part1-john-nirenberg,hanlin:weaksol-part2-local-boundedness}
\emph{Let $u \in H^1(\Omega)$ be a nonnegative supersolution in $\Omega$ in the following sense}
\[
(*) \qquad \int_\Omega a_{ij}D_i uD_j\varphi \ge \int_\Omega f\varphi \qquad \text{for any $\varphi \in H^1_0(\Omega)$ with $\varphi \ge 0$ in $\Omega$.}
\]
\emph{Suppose $f \in L^q(\Omega)$ for some $q > n/2$. Then for any $B_R \subset \Omega$, any $0 < p < n/(n-2)$ and any $0 < \theta < \tau < 1$}
\[
\inf_{B_{\theta R}} u + R^{2-\frac{n}{q}}\|f\|_{L^q(B_R)} \ge C\left(\frac{1}{R^n}\int_{B_{\tau R}} u^p\right)^{\frac{1}{p}},
\]
\emph{where $C$ is a positive constant depending only on $n$, $p$, $q$, $\lambda$, $\Lambda$, $\theta$ and $\tau$.}
\end{theorem}

% Source page 0089
\begin{proof}
We only prove for $R = 1$.

\noindent\textit{Step 1.} We prove that the result holds for some $p_0 > 0$.
Set $\baru = u + k > 0$, for some $k > 0$ to be determined and $v = \baru^{-1}$. First, we derive an equation for $v$. For any $\varphi \in H^1_0(B_1)$ with $\varphi \geq 0$ in $B_1$, consider $\baru^{-2}\varphi$ as the test function in (*). We have

\[
\int_{B_1} a_{ij}D_i u \frac{D_j\varphi}{\baru^2}
- 2\int_{B_1} a_{ij}D_i u\,D_j\baru \frac{\varphi}{\baru^3}
\geq \int_{B_1} f\frac{\varphi}{\baru^2}
\]

\noindent Note $D\baru = Du$ and $Dv = -\baru^{-2}D\baru$. Therefore, we obtain

\[
\int_{B_1} a_{ij}D_j v D_i\varphi + \tilde{f}\,v\varphi \leq 0,
\]

\noindent where

\[
\tilde{f} = \frac{f}{\baru}.
\]

\noindent In other words, $v$ is a nonnegative subsolution of some homogeneous equation. Choose $k = \|f\|_{L^q}$ if $f$ is not identical zero. Otherwise, choose an arbitrary $k > 0$ and then let $k \to 0+$. Note

\[
\|\tilde{f}\|_{L^q(B_1)} \leq 1.
\]

\noindent Theorem 4.1 implies that for any $\tau \in (\theta,1)$ and any $p > 0$

\[
\sup_{B_{\theta}} \baru^{-p} \leq C \int_{B_{\tau}} \baru^{-p},
\]

\noindent or,

\[
\inf_{B_{\theta}} \baru \geq C\left(\int_{B_{\tau}} \baru^{-p}\,dx\right)^{-1/p}
= C\left(\int_{B_{\tau}} \baru^{-p}\int_{B_{\tau}} \baru^p\right)^{-1/p}\left(\int_{B_{\tau}} \baru^p\right)^{1/p},
\]

\noindent where $C$ is a positive constant depending only on $n, q, p, \lambda, \Lambda, \tau$ and $\theta$.
The key point next is to show that there exists a $p_0 > 0$ such that

\[
\int_{B_{\tau}} \baru^{-p_0}\cdot \int_{B_{\tau}} \baru^{p_0} \leq C,
\]

\noindent where $C$ is a positive constant depending only on $n, q, \lambda, \Lambda$ and $\tau$. We will show for any $\tau < 1$

\begin{equation}
\int_{B_{\tau}} e^{p_0|w|} \leq C,
\tag{1}
\end{equation}

\noindent where $w = \log \baru - \beta$ with $\beta = |B_{\tau}|^{-1}\int_{B_{\tau}} \log \baru$.

\noindent We have two methods to proceed:

\noindent (i) To prove directly.

\noindent (ii) To prove that $w$ is BMO, i.e., for any $\Br{r}(y) \subset B_1$,

\[
\frac{1}{r^n}\int_{\Br{r}} |w-w_{\wy{r}}|\,dx \leq C.
\]

\noindent Then (1) follows from Theorem 3.5 in Chapter 3 (John-Nirenberg Lemma).

% Source page 0090

We first prove (1) directly. Recall $\bar u = u + k \ge k > 0$. Note that
\[
e^{p_0|w|} = 1 + p_0|w| + \frac{(p_0|w|)^2}{2!} + \cdots + \frac{(p_0|w|)^n}{n!} + \cdots .
\]
Hence we need to estimate
\[
\int_{B_\tau} |w|^\beta,
\]
for each positive integer $\beta$.

We first derive an equation for $w$. Consider $\bar u^{-1}\varphi$ as the test function in $(*)$.
Here we need $\varphi \in L^\infty(B_1)\cap H_0^1(B_1)$ with $\varphi \ge 0$. By a direct calculation as before
and by the fact $Dw = \bar u^{-1}Du$, we have
\begin{equation}
\int_{B_1} a_{ij} D_iw D_jw\varphi \le \int_{B_1} a_{ij} D_iw D_j\varphi + \int_{B_1} (-\tilde f\varphi)
\tag{2}
\end{equation}
for any $\varphi \in L^\infty(B_1)\cap H_0^1(B_1)$ with $\varphi \ge 0$.

Replace $\varphi$ by $\varphi^2$ in (2). Then the Cauchy inequality implies
\[
\int_{B_1} |Dw|^2\varphi^2 \le C\left\{\int_{B_1} |D\varphi|^2 + \int_{B_1} |\tilde f|\varphi^2\right\}.
\]
By the Hölder inequality and the Sobolev inequality, we obtain
\[
\int_{B_1} |\tilde f|\varphi^2 \le \|\tilde f\|_{L^{n/2}}\|\varphi\|_{L^{2n/(n-2)}}^2 \le c(n,q)\|D\varphi\|_{L^2}^2.
\]
Therefore, we have
\begin{equation}
\int_{B_1} |Dw|^2\varphi^2 \le C\int_{B_1} |D\varphi|^2,
\tag{3}
\end{equation}
where $C$ is a positive constant depending only on $n$, $q$, $\lambda$ and $\Lambda$. Take $\varphi \in C_0^1(B_1)$
with $\varphi \equiv 1$ in $B_\tau$. Then we obtain
\begin{equation}
\int_{B_\tau} |Dw|^2 \le C,
\tag{4}
\end{equation}
where $C$ is a positive constant depending only on $n$, $q$, $\lambda$, $\Lambda$ and $\tau$. Hence the
Poincaré inequality implies
\[
\int_{B_\tau} w^2 \le c(n,\tau)\int_{B_\tau} |Dw|^2 \le C,
\]
since $\int_{B_\tau} w = 0$. By (3), we conclude for any $\tau' \in (\tau,1)$
\begin{equation}
\int_{B_{\tau'}} w^2 \le C,
\tag{5}
\end{equation}
where $C$ is a positive constant depending only on $n$, $q$, $\lambda$, $\Lambda$, $\tau$ and $\tau'$.

Next, we estimate $\int_{B_{\tau'}} |w|^\beta$ for any $\beta \ge 2$. Choose $\varphi = \zeta^2 |w_m|^{2\beta} \in H_0^1(B_1)\cap
L^\infty(B_1)$ with
\[
w_m =
\begin{cases}
-m & w \le -m,\\
w & |w| < m,\\
m & w \ge m.
\end{cases}
\]
% Source page 0091

Substitute such a $\varphi$ in (2) to get

\[
\int_{B_1} \zeta^2|w_m|^{2\beta}a_{ij}D_iwD_jw
\le (2\beta)\int_{B_1} \zeta^2a_{ij}D_iwD_j|w_m||w_m|^{2\beta-1}
\]
\[
\qquad + \int_{B_1} 2\zeta|w_m|^{2\beta}a_{ij}D_iwD_j\zeta + \int_{B_1} |\tilde f|\zeta^2|w_m|^{2\beta}.
\]

Note

\[
a_{ij}D_iwD_j|w_m| = a_{ij}D_iw_mD_j|w_m| \le a_{ij}D_iw_mD_jw_m \quad \text{a.e. in } B_1.
\]

The Hölder inequality implies

\[
(2\beta)|w_m|^{2\beta-1} \le \frac{2\beta-1}{2\beta}|w_m|^{2\beta} + \frac{1}{2\beta}(2\beta)^{2\beta}
\]
\[
= \left(1-\frac{1}{2\beta}\right)|w_m|^{2\beta} + (2\beta)^{2\beta-1}.
\]

We obtain

\[
\int_{B_1} \zeta^2|w_m|^{2\beta}a_{ij}D_iwD_jw
\le \left(1-\frac{1}{2\beta}\right)\int_{B_1} \zeta^2|w_m|^{2\beta}a_{ij}D_iw_mD_jw_m
\]
\[
\qquad + (2\beta)^{2\beta-1}\int_{B_1}\zeta^2a_{ij}D_iw_mD_jw_m
\]
\[
\qquad + \int_{B_1} 2\zeta|w_m|^{2\beta}a_{ij}D_iwD_j\zeta + \int_{B_1} |\tilde f|\zeta^2|w_m|^{2\beta},
\]

and hence

\[
\int_{B_1} \zeta^2|w_m|^{2\beta}a_{ij}D_iwD_jw
\le (2\beta)^{2\beta}\int_{B_1}\zeta^2a_{ij}D_iw_mD_jw_m
\]
\[
\qquad + 4\beta\int_{B_1}\zeta|w_m|^{2\beta}a_{ij}D_iwD_j\zeta + 2\beta\int_{B_1} |\tilde f|\zeta^2|w_m|^{2\beta}.
\]

Therefore, we have

\[
\int_{B_1}\zeta^2|w_m|^{2\beta}|Dw|^2 \le C\left\{(2\beta)^{2\beta}\int_{B_1}\zeta^2|Dw_m|^2\right.
\]
\[
\qquad \left.+ \beta\int_{B_1}|w_m|^{2\beta}|Dw||D\zeta| + \beta\int_{B_1}|\tilde f|\zeta^2|w_m|^{2\beta}\right\}.
\]

Note that the first term in the right hand side is bounded in (4). Applying the
Cauchy inequality to the second term in the right hand side, we conclude

\[
\int_{B_1}\zeta^2|w_m|^{2\beta}|Dw|^2 \le C\left\{(2\beta)^{2\beta}\int_{B_1}\zeta^2|Dw_m|^2\right.
\]
\[
\qquad \left.+ \beta^2\int_{B_1}|w_m|^{2\beta}|D\zeta|^2 + \beta\int_{B_1}|\tilde f|\zeta^2|w_m|^{2\beta}\right\}.
\]

Note $Dw = Dw_m$ for $|w|<m$ and $Dw_m = 0$ for $|w|>m$. Hence we have

\[
\int_{B_1}\zeta^2|w_m|^{2\beta}|Dw_m|^2 \le C\left\{(2\beta)^{2\beta}\int_{B_1}\zeta^2|Dw_m|^2\right.
\]
\[
\qquad \left.+ \beta^2\int_{B_1}|w_m|^{2\beta}|D\zeta|^2 + \beta\int_{B_1}|\tilde f|\zeta^2|w_m|^{2\beta}\right\}.
\]
% Source page 0092

In the following, we write $w = w_m$ and then let $m \to +\infty$ at the end. By the
Hölder inequality, we obtain
\[
\lvert D(\zeta \lvert w^\beta \rvert)\rvert^2 \le 2 \lvert D\zeta \rvert^2 \lvert w \rvert^{2\beta}
+ 2\beta^2 \zeta^2 \lvert w \rvert^{2\beta-2} \lvert Dw \rvert^2
\]
\[
\le 2 \lvert D\zeta \rvert^2 \lvert w \rvert^{2\beta}
+ 2 \zeta^2 \lvert Dw \rvert^2 \left(\frac{\beta-1}{\beta} \lvert w \rvert^{2\beta}
+ \frac{1}{\beta}\beta^{2\beta}\right),
\]
and hence
\[
\int_{B_1} \lvert D(\zeta \lvert w^\beta \rvert)\rvert^2 \le C\left\{(2\beta)^{2\beta} \int_{B_1} \zeta^2 \lvert Dw \rvert^2
+ \beta^2 \int_{B_1} \lvert D\zeta \rvert^2 \lvert w \rvert^{2\beta}
+ \beta \int_{B_1} \lvert \tilde f \rvert \zeta^2 \lvert w \rvert^{2\beta}\right\}.
\]

Then the Hölder inequality implies
\[
\int_{B_1} \lvert \tilde f \rvert \zeta^2 \lvert w \rvert^{2\beta}
\le \left(\int_{B_1} \lvert \tilde f \rvert^q\right)^{\frac{1}{q}}
\left(\int_{B_1} (\zeta \lvert w^\beta \rvert)^{\frac{2q}{q-1}}\right)^{1-\frac{1}{q}}.
\]

Note $2^* = \frac{2n}{n-2} > \frac{2q}{q-1} > 2$ if $q > n/2$. By the interpolation inequality and the
Sobolev inequality, we have for any small $\varepsilon > 0$
\[
\lVert \zeta \lvert w^\beta \rvert \rVert_{L^{\frac{2q}{q-1}}}
\le \varepsilon \lVert \zeta \lvert w^\beta \rvert \rVert_{L^{2^*}}
+ C(n,q)\varepsilon^{-\frac{n}{2q-n}} \lVert \zeta \lvert w^\beta \rvert \rVert_{L^2}
\]
\[
\le \varepsilon \lVert D(\zeta \lvert w^\beta \rvert)\rVert_{L^2}
+ C(n,q)\varepsilon^{-\frac{n}{2q-n}} \lVert \zeta \lvert w^\beta \rvert \rVert_{L^2}.
\]

Therefore, we obtain by (3)
\[
\int_{B_1} \lvert D(\zeta \lvert w^\beta \rvert)\rvert^2 \le C\left\{(2\beta)^{2\beta} \int_{B_1} \zeta^2 \lvert Dw \rvert^2
+ \beta^\alpha \int_{B_1} \bigl(\lvert D\zeta \rvert^2 + \zeta^2\bigr)\lvert w \rvert^{2\beta}\right\}
\]
\[
\le C\left\{(2\beta)^{2\beta} \int_{B_1} \lvert D\zeta \rvert^2
+ \beta^\alpha \int_{B_1} \bigl(\lvert D\zeta \rvert^2 + \zeta^2\bigr)\lvert w \rvert^{2\beta}\right\},
\]
for some positive constant $\alpha$ depending only on $n$ and $q$. Apply the Sobolev inequality for
$\zeta \lvert w^\beta \rvert \in W_0^{1,2}(\mathbb{R}^n)$ with $\chi = \frac{n}{n-2}$ to get
\[
\left(\int_{B_1} \zeta^{2\chi}\lvert w \rvert^{2\beta\chi}\right)^{\frac{1}{\chi}}
\le C\beta^\alpha \left\{(2\beta)^{2\beta} \int_{B_1} \lvert D\zeta \rvert^2
+ \int_{B_1} \bigl(\lvert D\zeta \rvert^2 + \zeta^2\bigr)\lvert w \rvert^{2\beta}\right\}.
\]

Choose a cut-off function $\zeta$ as follows. For $\tau \le r < R \le 1$, set $\zeta \equiv 1$ on $B_r$, $\zeta \equiv 0$
in $B_1 \setminus B_R$ and $\lvert D\zeta \rvert \le \frac{2}{R-r}$. Therefore, we have
\[
\left(\int_{B_r} \lvert w \rvert^{2\beta\chi}\right)^{\frac{1}{\chi}}
\le \frac{C\beta^\alpha}{(R-r)^2}\left\{(2\beta)^{2\beta} + \int_{B_R} \lvert w \rvert^{2\beta}\right\}.
\]

For some $\tau' \in (\tau,1)$, set for any $i=1,2,\cdots$
\[
\beta_i = \chi^{i-1}, \qquad r_i = \tau + \frac{1}{2^{i-1}}(\tau' - \tau).
\]

Then for each $i=1,2,\cdots$, we have
\[
\left(\int_{B_{r_i}} \lvert w \rvert^{2\chi^i}\right)^{\frac{1}{\chi}}
\le \frac{C\chi^{(i-1)\alpha}2^{2(i-1)}}{(\tau' - \tau)^2}
\left\{(2\chi^{i-1})^{2\chi^{i-1}} + \int_{B_{r_{i-1}}} \lvert w \rvert^{2\chi^{i-1}}\right\}.
\]
% Source page 0093

Set

\[
I_j = \|w\|_{L^{2\chi^j}(B_{r_j})}.
\]

Then, we have for $j = 1, 2, \cdots$,
\[
I_j \le C^{\frac{j}{2\chi^j}}\{2\chi^{j-1} + I_{j-1}\},
\]
where $C$ is a positive constant depending only on $n$, $q$, $\lambda$, $\Lambda$, $\tau$ and $\tau'$. Iterating the above inequality and observing that
\[
\sum_{i=0}^{\infty} \frac{i}{\chi^i} < \infty,
\]
we obtain
\[
I_j \le C \sum_{i=1}^{j} \chi^{i-1} + CI_0,
\]
or,
\[
I_j \le C\chi^j + CI_0.
\]
Now for $\beta \ge 2$, there exists a $j$ such that $2\chi^{j-1} \le \beta < 2\chi^j$. Hence, we have
\[
I_{\beta}(B_{\tau}) \equiv \left(\int_{B_{\tau}} |w|^{\beta}\right)^{\frac{1}{\beta}} \le CI_j \le C\chi^j + CI_0
\]
\[
\le C\beta + CI_0 \le C_0\beta,
\]
since $I_0$ is bounded in (5). Then, we obtain for $\beta \ge 1$
\[
\int_{B_{\tau}} |w|^{\beta}\,dx \le C_0^{\beta}\beta^{\beta} \le C_0^{\beta} e^{\beta}\beta!.
\]
where we used the Sterling formula for the integer $\beta$. Hence, for any integer $\beta \ge 1$, we conclude
\[
\int_{B_{\tau}} \frac{(p_0|w|)^{\beta}}{\beta!} \le p_0^{\beta}(C_0e)^{\beta} \le \frac{1}{2^{\beta}},
\]
by choosing $p_0 = (2C_0e)^{-1}$. This proves that
\[
\int e^{p_0|w|} = \int 1 + p_0|w| + \frac{(p_0|w|)^2}{2!} + \cdots
\]
\[
\le 1 + \frac{1}{2^1} + \frac{1}{2^2} + \cdots \le 2.
\]

We remark that the above method is elementary in nature.
Now we give the second proof of the estimate (1) by using BMO. The estimate
(3) yields
\[
\int_{B_1} |Dw|^2 \zeta^2 \le C \int_{B_1} |D\zeta|^2 \qquad \text{for any } \zeta \in C_0^1(B_1).
\]
Then for any $B_{2r}(y) \subset B_1$, choose $\zeta$ with
\[
\operatorname{supp}\zeta \subset B_{2r}(y), \ \zeta \equiv 1 \text{ in } B_r(y), \ |D\zeta| \le \frac{2}{r}.
\]
We obtain
\[
\int_{B_r(y)} |Dw|^2 \le Cr^{n-2}.
\]
% Source page 0094

Hence the Poincar\'e inequality implies
\[
\frac{1}{r^n}\int_{B_r(y)} |w-w_{y,r}|
\leq \frac{1}{r^n}\left(\int_{B_r(y)} |w-w_{y,r}|^2\right)^{\frac12}
\leq \frac{1}{r^{\frac n2}}\left(r^2 \int_{B_r(y)} |Dw|^2\right)^{\frac12}
\leq C,
\]
or, $w \in \BMO$. Then the John-Nirenberg Lemma implies
\[
\int_{B_r} e^{p|w|} \leq C.
\]

\textit{Step 2.} The result holds for any positive $p<n/(n-2)$.\\
We need to prove, for any $0<r_1<r_2<1$ and $0<p_2<p_1<n/(n-2)$,
\begin{equation}
\left(\int_{B_{r_1}} u^{p_1}\right)^{\frac{1}{p_1}}
\leq C\left(\int_{B_{r_2}} u^{p_2}\right)^{\frac{1}{p_2}},
\tag{6}
\end{equation}
where $C$ is positive constant depending only on $n$, $q$, $\lambda$, $A$, $r_1$, $r_2$, $p_1$ and $p_2$. A
similar calculation may be found before. Here, we just point out some key steps.
Take $\varphi=\bar u^{-\beta}\eta^2$ for $\beta\in(0,1)$ as the test function in (*). Then we have
\[
\int_{B_1} |Du|^2 \bar u^{-\beta-1}\eta^2
\leq C\left\{\frac{1}{\beta^2}\int_{B_1} |D\eta|^2\bar u^{1-\beta}
+\frac{1}{\beta}\int_{B_1} \left|\frac{f}{k}\right| \eta^2 \bar u^{-\beta}\right\}.
\]
Set $\gamma=1-\beta\in(0,1)$ and $w=u^{\frac{\gamma}{2}}$. Then, we obtain
\[
\int |Dw|^2\eta^2 \leq \frac{C}{(1-\gamma)^{\alpha}} \int w^2(|D\eta|^2+\eta^2),
\]
or
\[
\int |D(w\eta)|^2 \leq \frac{C}{(1-\gamma)^{\alpha}} \int w^2(|D\eta|^2+\eta^2),
\]
for some positive $\alpha>0$. The Sobolev embedding theorem and an appropriate
choice of cut-off functions imply, with $\chi=n/(n-2)$, for any $0<r<R<1$
\[
\left(\int_{B_r} w^{2\chi}\right)^{\frac{1}{\chi}}
\leq \frac{C}{(1-\gamma)^{\alpha}}\frac{1}{(R-r)^2}\int_{B_R} w^2,
\]
or
\[
\left(\int_{B_r} \bar u^{\gamma\chi}\right)^{\frac{1}{\gamma\chi}}
\leq \left(\frac{C}{(1-\gamma)^{\alpha}}\frac{1}{(R-r)^2}\right)^{\frac{1}{\gamma}}
\left(\int_{B_R} \bar u^\gamma\right)^{\frac{1}{\gamma}}.
\]
This holds for any $\gamma\in(0,1)$. Now (6) follows after finitely many iterations.
\end{proof}

Now the Moser's Harnack inequality is an easy consequence of above results.

\begin{theorem}\label{hanlin:weaksol-part2-moser-harnack}\dcref{Theorem 4.16}\group{ch04-weak-solutions-part-ii}\source{han-lin-lecture-notes:page-0094,han-lin-lecture-notes:page-0095}\uses{hanlin:weaksol-part2-scaled-local-boundedness,hanlin:weaksol-part2-weak-harnack}
Let $u\in H^1(\Omega)$ be a nonnegative solution in $\Omega$
\[
\int_\Omega a_{ij}D_i uD_j\varphi = \int_\Omega f\varphi \qquad \text{for any }\varphi\in H^1_0(\Omega).
\]
Suppose $f\in L^q(\Omega)$ for some $q>n/2$. Then for any $B_R\subset \Omega$,
\[
\sup_{B_{R/2}} u \leq C\left\{ \inf_{B_{R/2}} u + R^{2-\frac{n}{q}}\|f\|_{L^q(B_R)} \right\}.
\]

% Source page 0095


where $C$ is a positive constant depending only on $n$, $\lambda$, $\Lambda$ and $q$.
\end{theorem}

The H\"older continuity follows easily from Theorem 4.16.

\begin{corollary}\label{hanlin:weaksol-part2-moser-holder}\dcref{Corollary 4.17}\group{ch04-weak-solutions-part-ii}\source{han-lin-lecture-notes:page-0095,han-lin-lecture-notes:page-0096}\uses{hanlin:weaksol-part2-scaled-local-boundedness,hanlin:weaksol-part2-moser-harnack,hanlin:weaksol-part2-oscillation-iteration}
Let $u \in H^1(\Omega)$ be a solution in $\Omega$, i.e.,
\[
\int_\Omega a_{ij} D_i u \, D_j \varphi = \int_\Omega f\varphi \qquad \text{for any } \varphi \in H_0^1(\Omega).
\]
Suppose $f \in L^q(\Omega)$ for some $q > n/2$. Then $u \in C^\alpha(\Omega)$ for some $\alpha \in (0,1)$ depending only on $n,q,\lambda$ and $\Lambda$. Moreover, for any $B_R \subset \Omega$
\[
|u(x)-u(y)| \le C \left( \frac{|x-y|}{R} \right)^\alpha
\left\{ \left( \frac{1}{R^n} \int_{B_R} u^2 \right)^{1/2} + R^{2-\frac{n}{q}} \|f\|_{L^q(B_R)} \right\}
\]
for any $x,y \in B_{R/2}$,

where $C$ is a positive constant depending only on $n$, $\lambda$, $\Lambda$ and $q$.
\end{corollary}

\begin{proof}
We prove the estimate for $R=1$. Set for $r \in (0,1)$
\[
M(r) = \sup_{B_r} u \qquad \text{and} \qquad m(r) = \inf_{B_r} u.
\]
Then $M(r) < +\infty$ and $m(r) > -\infty$. It suffices to show
\[
\omega(r) \equiv M(r)-m(r) \le Cr^\alpha \{\|u\|_{L^2(B_1)} + \|f\|_{L^q(B_1)}\} \qquad \text{for any } r < \frac12.
\]
Set $\delta = 2 - \frac{n}{q}$. Apply Theorem 4.16 to $M(r)-u \ge 0$ in $B_r$ to get
\[
\sup_{B_{r/2}} (M(r)-u) \le C \left\{ \inf_{B_{r/2}} (M(r)-u) + r^\delta \|f\|_{L^q(B_r)} \right\},
\]
i.e.,
\[
(1) \qquad M(r) - m\left( \frac r2 \right) \le C \left\{ M(r) - M\left( \frac r2 \right) + r^\delta \|f\|_{L^q(B_r)} \right\}.
\]
Similarly, apply Theorem 4.16 to $u - m(r) \ge 0$ in $B_r$ to get
\[
(2) \qquad M\left( \frac r2 \right) - m(r) \le C \left\{ m\left( \frac r2 \right) - m(r) + r^\delta \|f\|_{L^q(B_r)} \right\}.
\]
Then by adding (1) and (2) together, we obtain
\[
\omega(r) + \omega\left( \frac r2 \right) \le C \left\{ \omega(r) - \omega\left( \frac r2 \right) + r^\delta \|f\|_{L^q(B_r)} \right\},
\]
or
\[
\omega\left( \frac r2 \right) \le \gamma \omega(r) + Cr^\delta \|f\|_{L^q(B_r)},
\]
for some $\gamma = \frac{C-1}{C+1} < 1$.

Apply Lemma 4.18 below with $\mu$ chosen such that
\[
\alpha = (1-\mu)\log \gamma / \log \tau < \mu\delta.
\]
We obtain
\[
\omega(\rho) \le C\rho^\alpha \left\{ \omega\left( \frac12 \right) + \|f\|_{L^q(B_1)} \right\} \qquad \text{for any } \rho \in (0, \tfrac12].
\]
% Source page 0096

While Theorem 4.14 implies
\[
\omega\!\left(\frac{1}{2}\right)\le C\{\|u\|_{L^2(B_1)}+\|f\|_{L^q(B_1)}\}.
\]
This finishes the proof.
\end{proof}

\begin{lemma}\label{hanlin:weaksol-part2-oscillation-iteration}\dcref{Lemma 4.18}\group{ch04-weak-solutions-part-ii}\source{han-lin-lecture-notes:page-0096}
Let $\omega$ and $\sigma$ be non-decreasing functions in an interval $(0,R]$.
Suppose for some $0<\gamma,\tau<1$
\[
\omega(\tau r)\le \gamma\omega(r)+\sigma(r)\qquad\text{for any } r\le R.
\]
Then, for any $\mu\in(0,1)$ and $r\le R$,
\[
\omega(r)\le C\left\{\left(\frac{r}{R}\right)^{\alpha}\omega(R)+\sigma\!\left(r^{\mu}R^{1-\mu}\right)\right\},
\]
where $C$ is a positive constant depending only on $\gamma$, $\tau$ and $\alpha=(1-\mu)\log\gamma/\log\tau$.
\end{lemma}

\begin{proof}
Fix some $r_1\le R$. Then for any $r\le r_1$, we have
\[
\omega(\tau r)\le \gamma\omega(r)+\sigma(r_1),
\]
since $\sigma$ is nondecreasing. We now iterate this inequality to get for any positive integer $k$
\[
\omega(\tau^k r_1)\le \gamma^k\omega(r_1)+\sigma(r_1)\sum_{i=0}^{k-1}\gamma^i\le \gamma^k\omega(R)+\frac{\sigma(r_1)}{1-\gamma}.
\]
For any $r\le r_1$, we choose $k$ such that
\[
\tau^k r_1<r\le \tau^{k-1}r_1.
\]
Hence, we have
\[
\omega(r)\le \omega(\tau^{k-1}r_1)\le \gamma^{k-1}\omega(R)+\frac{\sigma(r_1)}{1-\gamma}
\]
\[
\le \frac{1}{\gamma}\left(\frac{r}{r_1}\right)^{\log\gamma/\log\tau}\omega(R)+\frac{\sigma(r_1)}{1-\gamma}.
\]
By letting $r_1=r^{\mu}R^{1-\mu}$, we obtain
\[
\omega(r)\le \frac{1}{\gamma}\left(\frac{r}{R}\right)^{(1-\mu)(\log\gamma/\log\tau)}\omega(R)+\frac{\sigma\!\left(r^{\mu}R^{1-\mu}\right)}{1-\gamma}.
\]
This finishes the proof.
\end{proof}

We also have the following Liouville theorem.

\begin{corollary}\label{hanlin:weaksol-part2-liouville}\dcref{Corollary 4.19}\group{ch04-weak-solutions-part-ii}\source{han-lin-lecture-notes:page-0096,han-lin-lecture-notes:page-0097}\uses{hanlin:weaksol-part2-moser-holder}
Suppose $u$ is a solution of a homogeneous equation in $\mathbb{R}^n$, i.e.,
\[
\int_{\mathbb{R}^n} a_{ij}D_i u D_j\varphi=0\qquad\text{for any }\varphi\in H_0^1(\mathbb{R}^n).
\]
If $u$ is bounded, then $u$ is a constant.
\end{corollary}

\begin{proof}
As in the proof of Corollary 4.17, we have for some $\gamma<1$
\[
\omega(r)\le \gamma\omega(2r)\qquad\text{for any } r>0.
\]
By an iteration, we obtain
\[
\omega(r)\le \gamma^k\omega(2^k r)\le C\gamma^k,
\]
% Source page 0097


since $u$ is bounded. Hence by letting $k \to \infty$, we conclude
\[
\omega(r)=0 \qquad \text{for any } r>0.
\]
Therefore, $u$ is constant.
\end{proof}

\section{Nonlinear Equations}

Up to now, we have been discussing linear equations of the following form
\[
-D_j(a_{ij}(x)D_i u)=f(x) \quad \text{in } \mathrm{B}_1.
\]
It is natural to ask how these results generalize to nonlinear equations. To answer
this question, let us consider an equation for a solution $v$ of the form
\[
v(x)=\Phi(u(x)),
\]
for some smooth function $\Phi:\R \to \R$ with $\Phi' \neq 0$. Any estimates for $u$ can be
translated to those for $v$. To find the equation for $v$, we write
\[
u=\Psi(v),
\]
with $\Psi=\Phi^{-1}$. Then by setting $\eta=\Psi'(v)\xi$ for $\xi \in \C_0^{\infty}(\mathrm{B}_1)$, we have
\[
\int a_{ij}D_i u D_j \xi
= \int a_{ij}\Psi'(v)D_i v D_j \xi
\]
\[
= \int a_{ij}D_i v D_j \eta - \int \frac{\Psi''(v)}{\Psi'(v)} a_{ij}D_i v D_j v\, \eta.
\]
Therefore if $u$ satisfies
\[
\int a_{ij}D_i u D_j \xi = \int f(x)\xi \qquad \text{for any } \xi \in \mathrm{H}_0^1(\mathrm{B}_1),
\]
then $v$ satisfies
\[
\int a_{ij}D_i v D_j \eta
= \int \left(\frac{\Psi''(v)}{\Psi'(v)}a_{ij}D_i v D_j v + \frac{1}{\Psi'(v)}f\right)\eta
\qquad \text{for any } \eta \in \C_0^{\infty}(\mathrm{B}_1).
\]
Note that the nonlinear term has a quadratic growth in terms of $Dv$. Hence, we
may extend the space of test functions to $\mathrm{H}_0^1(\mathrm{B}_1)\cap \Linf(\mathrm{B}_1)$. It turns out that
$\mathrm{H}^1(\mathrm{B}_1)\cap \Linf(\mathrm{B}_1)$ is also the right space for solutions. The following example
illustrates the boundedness of solutions is essential.

\begin{example}\label{hanlin:weaksol-part2-nonlinear-counterexample}\dcref{Example 4.20}\group{ch04-weak-solutions-part-ii}\source{han-lin-lecture-notes:page-0097}
Consider the equation
\[
-\Delta u = |Du|^2 \quad \text{in } B_R \subset \R^2,
\]
with $R<1$. It is easy to check that $u(x)=\log\log |x|^{-1}-\log\log R^{-1} \in \mathrm{H}^1(\mathrm{B}_R)$ is
a weak solution with zero boundary data. Note that $u(x)\equiv 0$ is also a solution.
\end{example}

In this section, we assume $a_{ij} \in \Linf(\mathrm{B}_1)$ satisfies
\[
\lambda |\xi|^2 \le a_{ij}(x)\xi_i\xi_j \le \Lambda |\xi|^2 \qquad \text{for any } x \in \mathrm{B}_1 \text{ and } \xi \in \R^n,
\]
for some positive constants $\lambda$ and $\Lambda$. We consider the nonlinear equation of the
form
\[
(\ast)\qquad \int a_{ij}(x)D_i u D_j \varphi = \int b(x,u,Du)\varphi
\qquad \text{for any } \varphi \in \mathrm{H}_0^1(\mathrm{B}_1)\cap \Linf(\mathrm{B}_1).
\]
% Source page 0098


We say the nonlinear term $b$ satisfies the natural growth condition if
\[
|b(x,u,p)| \le C(u)\bigl(f(x)+|p|^2\bigr) \qquad \text{for any } (x,u,p)\in B_1\times \R\times \R^n,
\]
for some constant $C(u)$ depending only on $u$, and $f\in L^q(B_1)$ for some $q\ge \frac{2n}{n+2}$.
We always assume
\[
u\in H^1(B_1)\cap L^\infty(B_1).
\]

\begin{lemma}\label{hanlin:weaksol-part2-nonlinear-harnack}\dcref{Lemma 4.21}\group{ch04-weak-solutions-part-ii}\source{han-lin-lecture-notes:page-0098,han-lin-lecture-notes:page-0099}\uses{hanlin:weaksol-part2-scaled-local-boundedness,hanlin:weaksol-part2-weak-harnack}
\textit{Suppose $u\in H^1(B_1)$ is a nonnegative solution of $(\ast)$ with $|u|\le M$ in $B_1$ and that $b$ satisfies the natural growth condition with $f(x)\in L^q(B_1)$ for some $q>\frac{n}{2}$. Then for any $B_R\subset B_1$}
\[
\sup_{B_{R/2}} u \le C\left\{\sup_{B_{R/2}} u + R^{2-\frac{n}{q}}\|f\|_{L^q(B_R)}\right\},
\]
\textit{where $C$ is a positive constant depending only on $n,\lambda,\Lambda,M$ and $q$.}
\end{lemma}

\begin{proof}
Let $v=\frac{1}{\alpha}(e^{\alpha u}-1)$ for some $\alpha>0$. Then for $\varphi\in H_0^1(B_1)\cap L^\infty(B_1)$ with $\varphi\ge 0$, we have
\[
\int a_{ij}D_i v D_j \varphi
= \int a_{ij}e^{\alpha u}D_i u D_j\varphi
\]
\[
= \int a_{ij}D_i u D_j(e^{\alpha u}\varphi)-\alpha \int a_{ij}e^{\alpha u}D_i u D_j u\,\varphi
\]
\[
= \int b(x,u,Du)e^{\alpha u}\varphi-\alpha \int a_{ij}e^{\alpha u}D_i u D_j u\,\varphi
\]
\[
\le C(M)\int \bigl(f(x)+|Du|^2\bigr)e^{\alpha u}\varphi-\alpha\lambda\int |Du|^2e^{\alpha u}\varphi.
\]
Hence by taking $\alpha$ large, we have
\[
\tag{1}
\int a_{ij}D_i v D_j\varphi \le C\int f(x)\varphi
\]
\[
\text{for any } \varphi\in H_0^1(B_1)\cap L^\infty(B_1)\text{ with }\varphi\ge 0,
\]
for some positive constant $C$ depending only on $n,\lambda,\Lambda$ and $M$. Observe that $u$ and $v$ are compatible. Therefore by Theorem 4.14, we obtain for any $p>0$
\[
\sup_{B_{R/2}} u \le C(M,\alpha)\sup_{B_{R/2}} v
\]
\[
\le C\left\{\left(\frac{1}{R^n}\int_{B_R} v^p\right)^{\frac{1}{p}} + R^{2-\frac{n}{q}}\left(\int_{B_R} f^q\right)^{\frac{1}{q}}\right\}
\]
\[
\le C\left\{\left(\frac{1}{R^n}\int_{B_R} u^p\right)^{\frac{1}{p}} + R^{2-\frac{n}{q}}\left(\int_{B_R} f^q\right)^{\frac{1}{q}}\right\}.
\]

For the lower bound, we let $w=\frac{1}{\alpha}(1-e^{-\alpha u})$. As before, by choosing $\alpha>0$ large we have
\[
\int a_{ij}D_i w D_j\varphi \ge C\int f(x)\varphi \qquad \text{for any } \varphi\in H_0^1(B_1)\cap L^\infty(B_1)\text{ with }\varphi\ge 0.
\]
% Source page 0099

Hence by Theorem 4.15, we obtain for any $p\in \left(0,\frac{n}{n-2}\right)$
\[
\left(\frac{1}{R^n}\int_{B_R} u^p\right)^{\frac{1}{p}}
\le C\left\{\inf_{B_{R/2}} u+R^{2-\frac{n}{q}}
\left(\int_{B_R} f^q\right)^{\frac{1}{q}}\right\}.
\]
We finish the proof by combining the above inequalities.
\end{proof}

\begin{remark}\label{hanlin:weaksol-part2-nonlinear-gradient-estimate}\dcref{Remark 4.22}\group{ch04-weak-solutions-part-ii}\source{han-lin-lecture-notes:page-0099}\uses{hanlin:weaksol-part2-nonlinear-harnack}
In estimate (1) in the above proof, take $\varphi=(u+M)\eta^2$
for some $\eta\in C_0^1(B_1)$. Then by the Hölder inequality, we conclude
\[
\int |Du|^2\eta^2 \le C\left\{\int \left(|D\eta|^2+|f|\eta^2\right)\right\},
\]
for some positive constant $C$ depending only on $n,\lambda,\Lambda$ and $M$. This implies the
interior $L^2$-estimate of the gradient $Du$ in terms of these constants together with
$\|f\|_{L^1(B_1)}$. This fact will be used in the proof of Theorem 4.24.
\end{remark}

\begin{corollary}\label{hanlin:weaksol-part2-nonlinear-holder}\dcref{Corollary 4.23}\group{ch04-weak-solutions-part-ii}\source{han-lin-lecture-notes:page-0099}\uses{hanlin:weaksol-part2-moser-holder,hanlin:weaksol-part2-nonlinear-harnack}
\emph{Suppose $u\in H^1(B_1)$ is a bounded solution of $(*)$ and that $b$
satisfies the natural growth condition with $f(x)\in L^q(B_1)$ for some $q>\frac{n}{2}$. Then $u\in
C^\alpha_{\mathrm{loc}}(B_1)$, for some $\alpha\in(0,1)$ depending only on $n,\lambda,\Lambda,q$
and $\|u\|_{L^\infty}$. Moreover,}
\[
|u(x)-u(y)|\le C|x-y|^\alpha \qquad \text{for any }x,y\in B_{1/2},
\]
\emph{where $C$ is a positive constant depending only on $n,\lambda,\Lambda,q,\|u\|_{L^\infty(B_1)}$
and $\|f\|_{L^q(B_1)}$.}
\end{corollary}

\begin{proof}
The proof is identical to that of Theorem 4.17, with Theorem 4.16
replaced by Lemma 4.21.
\end{proof}

\begin{theorem}\label{hanlin:weaksol-part2-nonlinear-gradient-holder}\dcref{Theorem 4.24}\group{ch04-weak-solutions-part-ii}\source{han-lin-lecture-notes:page-0099,han-lin-lecture-notes:page-0100}\uses{hanlin:weaksol-part1-campanato,hanlin:weaksol-part1-lemma-3-4,hanlin:weaksol-part1-harmonic-lemma,hanlin:weaksol-part1-grad-holder,hanlin:weaksol-part2-nonlinear-gradient-estimate,hanlin:weaksol-part2-nonlinear-holder}
\emph{Suppose $u\in H^1(B_1)$ is a bounded solution of $(*)$ and that $b$
satisfies the natural growth condition with $f\in L^q(B_1)$ for some $q>n$. Assume
further that $a_{ij}\in C^\alpha(B_1)$ for $\alpha=1-\frac{n}{q}$. Then $Du\in C^\alpha_{\mathrm{loc}}(B_1)$.
Moreover,}
\[
\|Du\|_{C^\alpha(B_{1/2})}\le C,
\]
\emph{where $C$ is a positive constant depending only on $n,\lambda,\Lambda,q,\|u\|_{L^\infty(B_1)}$
and $\|f\|_{L^q(B_1)}$.}
\end{theorem}

\begin{proof}
We only need to prove $Du\in L^\infty_{\mathrm{loc}}$. Then the Hölder continuity is
implied by Theorem 3.11 in Chapter 3. For any $B_r(x_0)\subset B_1$, solve for $w$ such that
\[
\int_{B_r(x_0)} a_{ij}(x_0)D_i w D_j\varphi = 0 \qquad \text{for any }\varphi\in H^1_0(B_r(x_0)),
\]
with $w-u\in H^1_0(B_r(x_0))$. Then the maximum principle implies
\[
\inf_{B_r(x_0)}u\le w\le \sup_{B_r(x_0)}u \qquad \text{in }B_r(x_0),
\]
or
\[
(1)\qquad \sup_{B_r(x_0)}|u-w|\le \osc_{B_r(x_0)}u.
\]
By Lemma 3.9 in Chapter 3, we have for any $0<\rho\le r$
\[
(2)\qquad \int_{B_\rho(x_0)} |Du|^2 \le c\left\{\left(\frac{\rho}{r}\right)^n \int_{B_r(x_0)} |Du|^2
+\int_{B_r(x_0)} |D(u-w)|^2\right\}.
\]
% Source page 0100



\[
\int_{B_\rho(x_0)} |Du-(Du)_{x_0,\rho}|^2 \le c\left\{\left(\frac{\rho}{r}\right)^{n+2}\int_{B_r(x_0)} |Du-(Du)_{x_0,r}|^2+\int_{B_r(x_0)} |D(u-w)|^2\right\}.
\tag{3}
\]

Note that the function $v=u-w\in H^1_0(B_r(x_0))$ satisfies
\[
\int_{B_r(x_0)} a_{ij}(x_0)D_ivD_j\varphi=\int_{B_r(x_0)} b(x,u,Du)\varphi
+\int_{B_r(x_0)} (a_{ij}(x_0)-a_{ij}(x))D_iuD_j\varphi
\]
for any $\varphi\in H^1_0(B_r(x_0))\cap L^\infty(B_r(x_0))$.

Taking $\varphi=v$ and by the Sobolev inequality, we obtain
\[
\int_{B_r(x_0)} |Dv|^2 \le C\left\{\int_{B_r(x_0)} |Du|^2|v|+r^{2\alpha}\int_{B_r(x_0)} |Du|^2+r^{n+2\alpha}\|f\|_{L^q(B_1)}^2\right\}.
\]
Hence with (1), we conclude
\[
\int_{B_r(x_0)} |Dv|^2 \le C\left\{\left(r^{2\alpha}+\osc_{B_r(x_0)}u\right)\int_{B_r(x_0)} |Du|^2+r^{n+2\alpha}\|f\|_{L^q}^2\right\}.
\tag{4}
\]

Corollary 4.23 implies $u\in C^{\delta_0}$ for some $\delta_0>0$. Therefore, we have by (2) and (4)
\[
\int_{B_\rho(x_0)} |Du|^2 \le C\left\{\left[\left(\frac{\rho}{r}\right)^n+r^{2\alpha}+r^{\delta_0}\right]\int_{B_r(x_0)} |Du|^2+r^{n+2\alpha}\|f\|_{L^q}^2\right\}.
\]

By Lemma 3.4 in Chapter 3, we obtain for any $\delta<1$ and any $B_r(x_0)\subset B_{7/8}$
\[
\int_{B_r(x_0)} |Du|^2 \le Cr^{n-2+2\delta}\left\{\int_{B_{7/8}} |Du|^2+\|f\|_{L^q(B_1)}^2\right\}.
\]

This implies $u\in C^\delta_{\mathrm{loc}}$ for any $\delta<1$. Moreover, for any $B_r(x_0)\subset B_{3/4}$
\[
\osc_{B_r(x_0)}u \le Cr^\delta,
\]
where $C$ is a positive constant depending only on $n,\lambda,\Lambda,q,\|u\|_{L^\infty(B_1)}$ and $\|f\|_{L^q(B_1)}$,
by the Remark 4.22. With (4), we have for any $B_r(x_0)\subset B_{2/3}$
\[
\int_{B_r(x_0)} |Dv|^2 \le C\left\{\left(r^{2\alpha}+r^\delta\right)r^{n-2+2\delta}\int_{B_{7/8}} |Du|^2+r^{n+2\alpha}\|f\|_{L^q}^2\right\}
\le Cr^{n+2\alpha'},
\]
for some $\alpha'<\alpha$ if $\delta\in(0,1)$ is chosen such that $3\delta>2$ and $\alpha+\delta>1$. Hence, with
(3) we obtain for any $B_r(x_0)\subset B_{2/3}$ and any $0<\rho\le r$
\[
\int_{B_\rho(x_0)} |Du-(Du)_{x_0,\rho}|^2 \le C\left\{\left(\frac{\rho}{r}\right)^{n+2}\int_{B_r(x_0)} |Du-(Du)_{x_0,r}|^2+r^{n+2\alpha'}\right\}.
\]

By Lemma 3.4 and Theorem 3.1 in Chapter 3 again, we conclude that $Du\in C^{\alpha'}_{\mathrm{loc}}$
for some $\alpha'<\alpha$, and in particular $Du\in L^\infty_{\mathrm{loc}}$. This finishes the proof.
\end{proof}
